\chapter{D-模的直像与逆像}

在本章中，$\varphi:X\to Y$ 总表示光滑代数簇之间的态射。我们以
$\mathcal{M}_L(\mathcal{D}_X)$（相应地 $\mathcal{M}_R(\mathcal{D}_X)$）记左（相应地右）
$\mathcal{D}_X$-模层的范畴，并略去下标时以 $D_X=\Gamma(X,\mathcal{D}_X)$ 记其整体微分算子环。

\section{双模 $\mathcal{D}_{X\to Y}$}

设 $X$ 是代数簇，$V$ 与 $W$ 是两个 $\mathcal{O}_X$-模。$V$ 到 $W$ 的 $k$-线性态射 $T$ 称为\emph{阶 $\le n$ 的微分态射}，如果对任意开集 $U$ 及任意 $(n+1)$-元组正则函数
$f_0,f_1,\dots,f_n\in R(U)$，在 $U$ 上有
$$[\dots[[T,f_0],f_1],\dots,f_n]=0.$$
以 $\Diff(V,W)$ 记 $V$ 到 $W$ 的所有微分态射构成的空间。同时规定：当 $p<0$ 时
$F_p\Diff(V,W)=0$，当 $p$ 为正时 $F_p\Diff(V,W)$ 为阶 $\le p$ 的微分态射构成的子空间。
显然 $F\Diff(V,W)$ 是 $\Diff(V,W)$ 穷尽的滤过。这个概念推广了 $X$ 上微分算子的概念：
若 $V=W=\mathcal{O}_X$，则 $\mathcal{O}_X$ 的微分自同态恰是 $X$ 上的微分算子。

类似地，可以定义 $V$ 到 $W$ 的微分态射层 $\mathcal{D}iff(V,W)$。

\begin{lemma}\label{lem:4.1.1}
设 $T$、$S$ 分别是阶 $\le n$、$\le m$ 的微分态射。则 $T\circ S$ 是阶 $\le n+m$ 的微分自同态。
\end{lemma}

\begin{proof}
对 $n+m$ 作归纳。若 $n=m=0$，则 $T$、$S$ 是 $\mathcal{O}_X$-模的态射，故 $T\circ S$ 也是
$\mathcal{O}_X$-模的态射，从而是阶 $\le 0$ 的微分态射。

现设 $n+m>0$。则
$$[T\circ S,f]=TSf-fTS=T[S,f]+[T,f]S,$$
且 $[T,f]$、$[S,f]$ 分别是阶 $\le n-1$ 与 $\le m-1$ 的微分态射。由归纳假设，这个微分态射的阶
$\le n+m-1$。因此 $T\circ S$ 的阶 $\le n+m$。
\end{proof}

因此，$\mathcal{O}_X$-模 $V$ 的所有微分自同态构成一个滤过环，而局部微分自同态构成一个滤过环层。

设 $X$、$Y$ 是光滑簇，$\varphi:X\to Y$ 是簇的态射。令
$\mathcal{D}_{X\to Y}=\varphi^*(\mathcal{D}_Y)$。则它既是 $\mathcal{O}_X$-模，又通过右乘成为右
$\varphi^{-1}\mathcal{D}_Y$-模。设 $\mathcal{C}$ 是 $\mathcal{O}_X$-模 $\mathcal{D}_{X\to Y}$ 的所有局部微分自同态、同时也是 $\varphi^{-1}\mathcal{D}_Y$-自同态的层。下面描述环层 $\mathcal{C}$ 的结构。

首先注意到 $\mathcal{O}_X$ 自然是 $F_0\mathcal{C}$ 的子环。

切层 $\mathcal{T}_Y$ 是 $\mathcal{D}_Y$ 的 $\mathcal{O}_Y$-子模。以 $\mathcal{J}_Y$ 记 $\mathcal{D}_Y$ 中由
$\mathcal{T}_Y$ 生成的左理想层。设 $y\in Y$，则存在 $y$ 的仿射开邻域 $U$ 及其上的坐标系
$(f_1,\dots,f_m;D_1,\dots,D_m)$，使得 $(D^I;I\in\mathbb{Z}_+^m)$ 是自由 $\mathcal{O}_U$-模 $\mathcal{D}_U$
关于左乘的一组基。左理想层 $\mathcal{J}_Y|_U$ 由 $(D^I;I\in\mathbb{Z}_+^m,\ |I|>0)$ 张成。因此有
$\mathcal{D}_Y=\mathcal{O}_Y\oplus\mathcal{J}_Y$。作为 $\mathcal{O}_X$-模，这导致直和分解
$$\mathcal{D}_{X\to Y}=\mathcal{O}_X\oplus\varphi^*(\mathcal{J}_Y).$$
以 $\alpha:\mathcal{D}_{X\to Y}\to\mathcal{O}_X$ 记相应的投影。

设 $S\in\mathcal{C}(U)$。我们断言 $\varphi^*(\mathcal{J}_Y)|_U$ 是 $S$-不变的。通过限制，可设 $U$ 是仿射的且
$\varphi(U)$ 包含在带有坐标系 $(f_1,\dots,f_m;D_1,\dots,D_m)$ 的“小”仿射开集 $V$ 中。于是
$\varphi^*(\mathcal{J}_Y)|_U$ 的整体截面是以 $(1\otimes D^I;I\in\mathbb{Z}_+^m,\ |I|>0)$ 为基的自由
$R(U)$-模。由于 $S$ 是右 $\mathcal{D}_V$-模 $\mathcal{D}_{X\to Y}(U)$ 的自同态，对任意 $g_I\in R(U)$ 有
$$S\Bigl(\sum_{|I|>0}g_I\otimes D^I\Bigr)=\sum_{|I|>0}S(g_I\otimes 1)D^I\in\varphi^*(\mathcal{J}_Y).$$
因此 $\varphi^*(\mathcal{J}_Y)|_U$ 在 $S$ 下不变。

由此可知 $\varphi^*(\mathcal{J}_Y)$ 是 $\mathcal{D}_{X\to Y}$ 的 $\mathcal{C}$-子模。这意味着商
$\mathcal{D}_{X\to Y}/\varphi^*(\mathcal{J}_Y)$ 是 $\mathcal{C}$-模。由上述讨论，自然单射
$\mathcal{O}_X\to\mathcal{D}_{X\to Y}$ 与该商映射的复合是 $\mathcal{O}_X$-模的同构。这个同构反过来定义了环层态射
$\gamma:\mathcal{C}\to\mathcal{E}nd_k(\mathcal{O}_X)$，由
$$\gamma_U(S)(f)=\alpha(S(f\otimes 1))$$
给出，其中 $f\in\mathcal{O}_X(U)$；显然 $\gamma$ 在 $\mathcal{O}_X$ 上为恒等。于是 $\gamma$ 把 $\mathcal{C}$ 映入
$\mathcal{O}_X$ 的微分自同态，即 $X$ 上的局部微分算子。因此构造出了滤过环层同态
$\gamma:\mathcal{C}\to\mathcal{D}_X$。本节的主要结果如下。

\begin{theorem}\label{thm:4.1.2}
映射 $\gamma:\mathcal{C}\to\mathcal{D}_X$ 是滤过环层的同构。
\end{theorem}

先证 $\gamma$ 是单射。从一个特殊情形开始。

\begin{lemma}\label{lem:4.1.3}
设 $U\subset X$ 是 $X$ 中的开集，$S$ 是 $\mathcal{C}(U)$ 中阶 $\le 0$ 的元素。则 $S=f\in\mathcal{O}_X(U)$。
\end{lemma}

\begin{proof}
由于 $S$ 的阶 $\le 0$，它实际上是 $\mathcal{O}_U$-模 $\mathcal{D}_{X\to Y}|_U$ 的自同态。只需证明 $S$ 在 $U$ 的一个开覆盖上的限制由函数给出即可。因此可设 $U$ 是仿射的且 $\varphi(U)$ 包含在 $Y$ 中带坐标系
$(f_1,\dots,f_m;D_1,\dots,D_m)$ 的“小”仿射开集 $V$ 内。更确切地说，由 2.10 与 3.5，可设 $Y$ 上存在坐标系
$(f_1,\dots,f_m;D_1,\dots,D_m)$ 使得
\begin{enumerate}[label=(\roman*)]
  \item $[D_i,D_j]=0$，$1\le i\le m$；
  \item $D^I$（$I\in\mathbb{Z}_+^m$）是自由 $R(V)$-模 $\mathcal{D}_V$（关于左乘）的一组基。
\end{enumerate}
于是在此情形下 $\mathcal{D}_{X\to Y}(U)=R(U)\otimes_{R(V)}\mathcal{D}_V$ 是以 $1\otimes D^I$（$I\in\mathbb{Z}_+^m$）为基的自由 $R(U)$-模。由于 $S$ 与 $R(X)$ 元素的左乘及 $\mathcal{D}_Y$ 元素的右乘都可交换，它完全由其在 $1\otimes 1$ 上的值决定。设
$S(1\otimes 1)=\sum a_I\otimes D^I$，其中 $a_I\in R(X)$。则
$$\sum a_I\otimes D^I f_j=S(1\otimes 1)f_j=S(1\otimes f_j)=S(f_j\circ\varphi\otimes 1)
=\sum a_I(f_j\circ\varphi)\otimes D^I=\sum a_I\otimes f_jD^I$$
对任意 $1\le j\le m$ 成立。这意味着
$$\sum a_I\otimes[D^I,f_j]=0$$
对任意 $1\le j\le m$ 成立。设 $I=(i_1,\dots,i_m)$。显然若 $i_j=0$ 则 $[D^I,f_j]=0$；若 $i_j>0$，令
$I'=(i_1,\dots,i_{j-1},i_j-1,i_{j+1},\dots,i_m)$，则 $[D^I,f_j]=i_jD^{I'}$。这立即使对所有 $I\ne 0$ 有 $a_I=0$。因此 $S=a_0$。
\end{proof}

现在证明 $\gamma_U$ 的单射性。设 $S\in\mathcal{C}(U)$ 且 $\gamma_U(S)=0$。对 $S$ 的阶 $p$ 作归纳证明 $S=0$。若 $p=0$，由引理 \ref{lem:4.1.3} 有 $S=f\in\mathcal{O}_X(U)$ 且 $\gamma_U(S)=\gamma_U(f)=f$，故 $S=0$。设结论对所有阶 $\le p-1$（$p>1$）的 $T\in\mathcal{C}(V)$ 及所有开集 $V\subset X$ 成立。若 $S$ 的阶 $\le p$，则对任意 $g\in\mathcal{O}_X(V)$、$V\subset U$，有
$\gamma([S,g])=[\gamma(S),g]=0$，而 $[S,g]$ 的阶 $\le p-1$，由归纳假设 $[S,g]=0$。这意味着 $S$ 的阶 $\le 0$，由证明的第一部分得 $S=0$。

这表明 $\gamma:\mathcal{C}\to\mathcal{D}_X$ 是单射。还需证明 $\gamma$ 是满射。

先证明所有局部向量场都在 $\gamma$ 的像中。以 $\mathcal{T}_X$、$\mathcal{T}_Y$ 分别记 $X$、$Y$ 的切层。由 2.12，$\mathcal{O}_Y$-模 $\mathcal{T}_Y$ 是有限秩局部自由的，因此其逆像 $\varphi^*(\mathcal{T}_Y)$ 是有限秩局部自由 $\mathcal{O}_X$-模。于是 $\varphi^*(\mathcal{T}_Y)$ 在开集 $U\subset X$ 上的截面完全由它在所有 $x\in U$ 的几何纤维中的像决定：
$$T_x(\varphi^*(\mathcal{T}_Y))=\mathcal{O}_{X,x}/\mathfrak{m}_{X,x}\otimes_{\mathcal{O}_{Y,\varphi(x)}}\mathcal{T}_{Y,\varphi(x)}=T_{\varphi(x)}(Y).$$
设 $\Psi\in\operatorname{Hom}_{\mathcal{O}_X}(\mathcal{T}_X,\varphi^*(\mathcal{T}_Y))$。则对任意 $x\in X$，这个态射决定几何纤维之间的线性映射
$$T_x(\Psi):T_x(X)=T_x(\mathcal{T}_X)\longrightarrow T_x(\varphi^*(\mathcal{T}_Y))=T_{\varphi(x)}(Y).$$

\begin{lemma}\label{lem:4.1.4}
存在唯一的 $\mathcal{O}_X$-模态射 $\Phi:\mathcal{T}_X\to\varphi^*(\mathcal{T}_Y)$，使得对任意 $x\in X$，由它在几何纤维上诱导的线性映射 $T_x(\Phi)$ 等于切线性映射
$T_x(\varphi):T_x(X)\to T_{\varphi(x)}(Y)$。
\end{lemma}

\begin{proof}
唯一性由上注明显。设 $U$ 是 $X$ 中开集，$T$ 是 $U$ 上向量场。设 $\mathcal{U}=(U_i;1\le i\le n)$ 是 $U$ 的一个开覆盖，其元素的像包含在 $Y$ 的“小”开集内，即 $\varphi(U_i)\subset V_i$，且在 $V_i$ 上存在坐标系
$(f_1,\dots,f_m;D_1,\dots,D_m)$。此时对任意 $x\in U_i$ 的切向量 $\xi\in T_x(X)$，有
$T_x(\varphi)(\xi)=\sum_{j=1}^m a_jD_j(\varphi(x))$。此外，对任意 $1\le j\le m$，
$$\bigl(T_x(\varphi)(\xi)\bigr)(df_j(\varphi(x)))=\sum_{k=1}^m a_kD_k(\varphi(x))(df_j(\varphi(x)))
=\sum_{k=1}^m a_kD_k(f_j)(\varphi(x))=a_j.$$
因此
$$a_j=\xi(T_x^*(\varphi)(df_j(\varphi(x))))=\xi(d(f_j\circ\varphi)(x)).$$
这最终给出
$$T_x(\varphi)(\xi)=\sum_{j=1}^m\xi(d(f_j\circ\varphi)(x))D_j(\varphi(x)).$$
函数 $x\mapsto\bigl((T|_{U_i})(f_j\circ\varphi)\bigr)(x)$（$1\le j\le m$）在 $U_i$ 上正则。于是
$$\sum_{j=1}^m(T|_{U_i})(f_j\circ\varphi)\otimes D_j$$
是 $\varphi^*(\mathcal{T}_Y)$ 在 $U_i$ 上的一个截面。此外，由上述公式，对任意 $x\in U_i$，这个截面在 $\varphi^*(\mathcal{T}_Y)$ 于 $\varphi(x)$ 的几何纤维 $T_{\varphi(x)}(Y)$ 中的像等于
$$\sum_{j=1}^m T(f_j\circ\varphi)(x)D_j(\varphi(x))=\sum_{j=1}^m T(x)(d(f_j\circ\varphi)(x))D_j(\varphi(x))=T_x(\varphi)(T(x)).$$
因此在 $U_i$ 上构造了一个截面，其在 $x\in U_i$ 的几何纤维 $T_{\varphi(x)}(Y)$ 中的像为 $T_x(\varphi)(T(x))$。由于截面完全由其在几何纤维中的像决定，这个截面唯一。可将这些截面在 $\mathcal{U}$ 的所有元素上粘合，得到 $U$ 上的截面 $\Phi(T)$。显然映射 $\Phi:T\mapsto\Phi(T)$ 具有所需性质。
\end{proof}

现在注意到局部地 $\mathcal{T}_Y$ 是 $\mathcal{D}_Y$ 的 $\mathcal{O}_Y$-模直和因子，故自然态射
$\varphi^*(\mathcal{T}_Y)\to\mathcal{D}_{X\to Y}$ 是单射，可将 $\varphi^*(\mathcal{T}_Y)$ 与 $\mathcal{D}_{X\to Y}$ 的子模等同。因此可把 $\Phi$ 看作 $\mathcal{T}_X$ 到 $\mathcal{D}_{X\to Y}$ 的态射。

\begin{lemma}\label{lem:4.1.5}
设 $U\subset X$ 是开集，$T$ 是 $U$ 上向量场。则存在唯一元素 $\delta_U(T)\in\mathcal{C}(U)$ 使得
$$\delta_U(T)|_V(g\otimes 1)=T(g)\otimes 1+g\Phi(T)|_V$$
对任意 $g\in\mathcal{O}_X(V)$ 及任意开集 $V\subset U$ 成立。
\end{lemma}

\begin{proof}
假设存在 $\mathcal{C}(U)$ 的元素 $S$ 满足
$$S|_V(g\otimes 1)=T(g)\otimes 1+g\Phi(T)|_V$$
对任意 $g\in\mathcal{O}_X(V)$ 及任意开集 $V\subset U$ 成立。设 $x\in U$，则存在 $\varphi(x)$ 在 $Y$ 中的仿射开邻域 $W$ 及 $x$ 的仿射开邻域 $V\subset U$ 使 $\varphi(V)\subset W$。于是 $\mathcal{D}_{X\to Y}(V)=R(V)\otimes_{R(W)}\mathcal{D}_W$。对任意 $g\in R(V)$、$Q\in\mathcal{D}_W$，由于 $S|_V$ 是右 $\mathcal{D}_W$-模自同态，有 $S|_V(g\otimes Q)=S(g\otimes 1)Q$。由此可知这样的 $S|_V$ 是唯一的，故 $S$ 由上述性质唯一确定。

由唯一性，为证存在性，只需对 $U$ 换成其开覆盖的元素来证明。因此一开始就可设 $U$ 是仿射开集且 $\varphi(U)$ 包含在 $Y$ 中带坐标系 $(f_1,\dots,f_m;D_1,\dots,D_m)$ 的“小”仿射开集 $W$ 内。如上，此时 $\mathcal{D}_{X\to Y}(V)=R(V)\otimes_{R(W)}\mathcal{D}_W$。首先定义双线性映射 $R(V)\times\mathcal{D}_W\to R(V)\otimes_{R(W)}\mathcal{D}_W$ 为
$$(g,S)\longmapsto T(g)\otimes S+g\Phi(T)|_V S.$$
现取 $h\in R(W)$，其与 $\varphi|_V$ 的复合是 $V$ 上正则函数。我们断言
$$\Phi(T)|_V h-(h\circ\varphi)\Phi(T)|_V=T(h\circ\varphi)\otimes 1.$$
利用引理 \ref{lem:4.1.4} 证明中的计算，有
$$\Phi(T)|_V h-(h\circ\varphi)\Phi(T)|_V
=\sum_{j=1}^m\bigl(T(f_j\circ\varphi)\otimes D_jh-(h\circ\varphi)T(f_j\circ\varphi)\otimes D_j\bigr)$$
$$=\sum_{j=1}^m T(f_j\circ\varphi)\otimes[D_j,h]
=\sum_{j=1}^m T(f_j\circ\varphi)\otimes D_j(h)
=\sum_{j=1}^m T(f_j\circ\varphi)(D_j(h)\circ\varphi)\otimes 1,$$
即这个表达式是一个函数。它在 $x\in V$ 处的值为
$$\sum_{j=1}^m T(f_j\circ\varphi)(x)D_j(h)(\varphi(x))
=\sum_{j=1}^m T(x)(d(f_j\circ\varphi)(x))D_j(\varphi(x))(dh(\varphi(x)))$$
$$=\bigl(T_x(\varphi)(T(x))\bigr)(dh(\varphi(x)))=T(x)(d(h\circ\varphi)(x))=T(h\circ\varphi)(x),$$
这证明了上述关系。因此对任意 $h\in R(Y)$ 有
\begin{align*}
T(g(h\circ\varphi))\otimes S+g(h\circ\varphi)\Phi(T)S
&=T(g)\otimes hS+gT(h\circ\varphi)\otimes S+g(h\circ\varphi)\Phi(T)S\\
&=T(g)\otimes hS+g\Phi(T)hS,
\end{align*}
故该映射可分解通过 $\mathcal{D}_{X\to Y}(V)$，从而定义了 $\delta_V(T)\in\mathcal{E}nd_k(\mathcal{D}_{X\to Y}(V))$。若 $V_f$（$f\in R(V)$）是 $V$ 中主开集，则
$\mathcal{D}_{X\to Y}(V_f)=R(V_f)\otimes_{R(W)}\mathcal{D}_W=R(V)_f\otimes_{R(W)}\mathcal{D}_W$，因此由局部化，$\delta_V(T)$“延拓”为 $\mathcal{D}_{X\to Y}(V_f)$ 的自同态，它本质上由同一公式给出，因而等于 $\delta_{V_f}(T)$。这意味着 $\delta_V(T)$ 定义了 $\mathcal{E}nd_k(\mathcal{D}_{X\to Y})$ 的元素。显然对任意 $g,h\in R(V)$，
$$[\delta_V(T),g](h\otimes 1)=\delta_V(T)(gh\otimes 1)-g\delta_V(T)(h\otimes 1)
=T(gh)\otimes 1-gT(h)\otimes 1=T(g)(h\otimes 1),$$
故 $\delta_V(T)$ 是微分自同态。此外它还与 $\mathcal{D}_W$ 的右作用可交换，故 $\delta_V(T)\in\mathcal{C}(V)$。存在性得证。
\end{proof}

另一方面，由于 $\varphi^*(\mathcal{T}_Y)\subset\varphi^*(\mathcal{J}_Y)=\ker\alpha$，对任意 $g\in\mathcal{O}_X$ 有
$$\gamma(\delta(T))(g)=\alpha(\delta(T)(g\otimes 1))=\alpha(T(g)\otimes 1)=T(g),$$
即 $(\gamma\circ\delta)(T)=T$。这说明 $\gamma$ 的像包含 $X$ 上所有向量场，而由 3.7.(ii)，这些向量场与 $\mathcal{O}_X$ 生成环层 $\mathcal{D}_X$，故 $\gamma$ 是满射。定理 \ref{thm:4.1.2} 得证。

因此可把 $\mathcal{D}_{X\to Y}$ 视为双模层，带左 $\mathcal{D}_X$-作用与右 $\varphi^{-1}(\mathcal{D}_Y)$-作用。

现在讨论 $X=k^n$、$Y=k^m$ 的情形。此时在 I.12 中构造了 $\mathcal{D}_{X\to Y}$ 上由
$$\frac{\partial}{\partial x_i}(P\otimes S)=\frac{\partial P}{\partial x_i}\otimes S
+\sum_{j=1}^m P\frac{\partial(y_j\circ\varphi)}{\partial x_i}\otimes\frac{\partial}{\partial y_j}S$$
给出的左 $\mathcal{D}_X$-模结构，其中 $P\in R(X)$、$S\in\mathcal{D}_Y$。显然这个左 $\mathcal{D}_X$-作用与刚才一般定义的作用一致。因此双模层 $\mathcal{D}_{X\to Y}$ 是 I 中多项式映射情形构造的自然推广。

\begin{remark}\label{rem:4.1.6}
设 $U$ 是光滑代数簇 $X$ 的开子集，考虑自然包含 $i:U\to X$。则
$i^{-1}(\mathcal{D}_X)=\mathcal{D}_X|_U=\mathcal{D}_U$。因此 $\mathcal{D}_{U\to X}$ 作为右 $i^{-1}(\mathcal{D}_U)$-模同构于 $\mathcal{D}_U$。显然 $\mathcal{D}_{U\to X}$ 通过左乘具有左 $\mathcal{D}_U$-模结构，此作用与右 $i^{-1}(\mathcal{D}_X)$-作用可交换。于是得到一个自然态射 $\mathcal{D}_U\to\mathcal{C}$。由 $\gamma:\mathcal{C}\to\mathcal{D}_U$ 的定义知，对任意开集 $V\subset U$ 与向量场 $T\in\mathcal{D}_U(V)$，
$$\gamma_V(T)(f)=\alpha(T(f))=\alpha([T,f]+fT)=T(f).$$
故自然包含与 $\gamma$ 的复合在局部向量场上是同构。由于 $\mathcal{D}_U$ 局部由向量场生成，这说明该复合是 $\mathcal{D}_U$ 的同构。因此 $\mathcal{D}_U$ 在双模 $\mathcal{D}_{U\to X}$ 中的作用就是 $\mathcal{D}_U$ 的自然作用。
\end{remark}

设 $X$、$Y$、$Z$ 是光滑簇，$\varphi:X\to Y$、$\psi:Y\to Z$ 是簇的态射。则
$$\mathcal{D}_{X\to Z}=(\psi\circ\varphi)^*(\mathcal{D}_Z)=\varphi^*(\psi^*(\mathcal{D}_Z))
=\varphi^*(\mathcal{D}_{Y\to Z})$$
$$=\mathcal{O}_X\otimes_{\varphi^{-1}(\mathcal{O}_Y)}\varphi^{-1}(\mathcal{D}_{Y\to Z})
=\mathcal{O}_X\otimes_{\varphi^{-1}(\mathcal{O}_Y)}\bigl(\varphi^{-1}(\mathcal{D}_Y)
\otimes_{\varphi^{-1}(\mathcal{D}_Y)}\varphi^{-1}(\mathcal{D}_{Y\to Z})\bigr)$$
$$=\bigl(\mathcal{O}_X\otimes_{\varphi^{-1}(\mathcal{O}_Y)}\varphi^{-1}(\mathcal{D}_Y)\bigr)
\otimes_{\varphi^{-1}(\mathcal{D}_Y)}\varphi^{-1}(\mathcal{D}_{Y\to Z})
=\mathcal{D}_{X\to Y}\otimes_{\varphi^{-1}(\mathcal{D}_Y)}\varphi^{-1}(\mathcal{D}_{Y\to Z}).$$
这个同构显然与 $\mathcal{D}_{X\to Z}$ 和 $\mathcal{D}_{X\to Y}$ 中由左乘给出的 $\mathcal{O}_X$-模结构相容。对于由右乘给出的右
$(\psi\circ\varphi)^{-1}(\mathcal{D}_Z)=\varphi^{-1}(\psi^{-1}(\mathcal{D}_Z))$-模结构，以及右
$\varphi^{-1}(\mathcal{D}_{Y\to Z})$-模结构，同样成立。我们断言左 $\mathcal{D}_X$-模作用也相容。

设 $x\in X$，$T$ 是 $x$ 的仿射开邻域 $U$ 上的向量场。则
$$(\mathcal{D}_{X\to Y}\otimes_{\varphi^{-1}(\mathcal{D}_Y)}\varphi^{-1}(\mathcal{D}_{Y\to Z}))_x
=\mathcal{D}_{X\to Y,x}\otimes_{\mathcal{D}_{Y,\varphi(x)}}\mathcal{D}_{Y\to Z,\varphi(x)}.$$
设 $f$ 是 $x$ 处正则函数的芽。则在 $\mathcal{D}_{X\to Y}$ 中
$$T(f\otimes 1)=T(f)\otimes 1+\sum_i g_i\otimes T_i,$$
其中 $g_i\in\mathcal{O}_{X,x}$、$T_i\in\mathcal{J}_{Y,\varphi(x)}$。这意味着 $T$ 在茎
$\mathcal{D}_{X\to Y,x}\otimes_{\mathcal{D}_{Y,\varphi(x)}}\mathcal{D}_{Y\to Z,\varphi(x)}$ 的元素
$(f\otimes 1)\otimes(1\otimes 1)$ 上的作用为
$$T((f\otimes 1)\otimes(1\otimes 1))=T(f\otimes 1)\otimes(1\otimes 1)$$
$$=(T(f)\otimes 1)\otimes(1\otimes 1)+\sum_i(g_i\otimes T_i)\otimes(1\otimes 1)
=(T(f)\otimes 1)\otimes(1\otimes 1)+\sum_i(g_i\otimes 1)\otimes T_i(1\otimes 1).$$
前面已指出 $\psi^*(\mathcal{J}_Z)_{\varphi(x)}$ 是 $\mathcal{D}_{Y,\varphi(x)}$-不变的。此外，由作用的构造，$\varphi(x)$ 处任一局部向量场 $S$ 把 $1\otimes 1$ 映入 $\psi^*(\mathcal{J}_Z)_{\varphi(x)}$。这说明对所有 $i$ 有 $T_i(1\otimes 1)\in\psi^*(\mathcal{J}_Z)_{\varphi(x)}$。因此在上述同构下，$T((f\otimes 1)\otimes(1\otimes 1))$ 映为
$T(f)\otimes 1+\sum_j h_j\otimes S_j$，其中 $S_j\in\mathcal{J}_{Z,\psi(\varphi(x))}$。这说明 $T$ 的两个作用一致。因此建立了如下结果。

\begin{proposition}\label{prop:4.1.7}
设 $X$、$Y$、$Z$ 是光滑代数簇，$\varphi:X\to Y$、$\psi:Y\to Z$ 是代数簇的态射。则作为双模层有
$$\mathcal{D}_{X\to Z}=\mathcal{D}_{X\to Y}\otimes_{\varphi^{-1}(\mathcal{D}_Y)}\varphi^{-1}(\mathcal{D}_{Y\to Z}).$$
\end{proposition}
\begin{annote}
$f_+,f^+$ 与 $\boxtimes$、$\otimes^L$、$R\mathcal H om$、$R\Gamma_Z$ 合起来是 $\mathcal D$-模的六种运算。$f_+$ 有两套等价定义：经双边模 $\mathcal D_{Y\leftarrow X}=\mathcal D_Y\otimes_{f^{-1}\mathcal D_X}f^{-1}\mathcal D_X$ 的张量，或经相对 $\mathcal D_{X\to Y}$ 的转移；$f^+$ 是 $f_+$ 的右伴随、只能作导出函子。整套形式对应 nLab 的 six operations 公理，其中基变换为 Beck–Chevalley：$(p_2)_!\circ(p_1)^*\cong g^*\circ f_!$。[来源：nLab 条目 six+operations（Last revised 2026-09-30）；原书 §1–§2, p97–108]
\end{annote}



\section{仿射簇的逆像与直像}

设 $X$、$Y$ 是两个光滑仿射簇，$\varphi:X\to Y$ 是代数簇的态射。则
$D_{X\to Y}=\Gamma(X,\mathcal{D}_{X\to Y})$ 具有自然的（左 $D_X$、右 $D_Y$）-双模结构。定义以下两个右正合函子：

\begin{enumerate}[label=(\roman*)]
  \item 逆像函子 $\varphi^+:\mathcal{M}_L(D_Y)\to\mathcal{M}_L(D_X)$，对任意左 $D_Y$-模 $M$ 令
    $$\varphi^+(M)=D_{X\to Y}\otimes_{D_Y}M;$$
  \item 直像函子 $\varphi_+:\mathcal{M}_R(D_X)\to\mathcal{M}_R(D_Y)$，对任意右 $D_X$-模 $N$ 令
    $$\varphi_+(N)=N\otimes_{D_X}D_{X\to Y}.$$
\end{enumerate}
显然这些定义推广了 I.12 中仿射空间情形的定义。

设 $Z$ 是另一个光滑仿射簇，$\psi:Y\to Z$ 是簇的态射。则由前述结果，$D_{X\to Y}\otimes_{D_Y}D_{Y\to Z}$ 到 $D_{X\to Z}$ 有自然态射，且与左 $D_X$-模和右 $D_Z$-模作用相容。

\begin{proposition}\label{prop:4.2.1}
$D_{X\to Z}=D_{X\to Y}\otimes_{D_Y}D_{Y\to Z}$。
\end{proposition}

\begin{proof}
由前述结果，
$$D_{X\to Z}=\Gamma(X,\mathcal{D}_{X\to Z})=\Gamma(X,(\psi\circ\varphi)^*(\mathcal{D}_Z))
=\Gamma(X,\varphi^*(\psi^*(\mathcal{D}_Z)))$$
$$=\Gamma(X,\varphi^*(\mathcal{D}_{Y\to Z}))=R(X)\otimes_{R(Y)}\Gamma(Y,\mathcal{D}_{Y\to Z})
=R(X)\otimes_{R(Y)}D_{Y\to Z}$$
$$=R(X)\otimes_{R(Y)}(D_Y\otimes_{D_Y}D_{Y\to Z})
=(R(X)\otimes_{R(Y)}D_Y)\otimes_{D_Y}D_{Y\to Z}=D_{X\to Y}\otimes_{D_Y}D_{Y\to Z}.$$
\end{proof}

这立即蕴含：对任意左 $D_Z$-模 $M$，
$$(\varphi^+\circ\psi^+)(M)=\varphi^+(\psi^+(M))=D_{X\to Y}\otimes_{D_Y}(D_{Y\to Z}\otimes_{D_Z}M)$$
$$=(D_{X\to Y}\otimes_{D_Y}D_{Y\to Z})\otimes_{D_Z}M=D_{X\to Z}\otimes_{D_Z}M=(\psi\circ\varphi)^+(M);$$
对任意右 $D_X$-模 $N$，
$$(\psi_+\circ\varphi_+)(N)=\psi_+(\varphi_+(N))=(N\otimes_{D_X}D_{X\to Y})\otimes_{D_Y}D_{Y\to Z}$$
$$=N\otimes_{D_X}(D_{X\to Y}\otimes_{D_Y}D_{Y\to Z})=N\otimes_{D_X}D_{X\to Z}=(\psi\circ\varphi)_+(N).$$
因此建立了如下结果。

\begin{theorem}\label{thm:4.2.2}
设 $X$、$Y$、$Z$ 是光滑仿射簇，$\varphi:X\to Y$、$\psi:Y\to Z$ 是簇的态射。则
\begin{enumerate}[label=(\roman*)]
  \item $\varphi^+\circ\psi^+=(\psi\circ\varphi)^+$；
  \item $\psi_+\circ\varphi_+=(\psi\circ\varphi)_+$。
\end{enumerate}
\end{theorem}

现研究逆像 $\varphi^+$ 与直像 $\varphi_+$ 的左导出函子。先从以下结果开始。

\begin{lemma}\label{lem:4.2.3}
设 $V$ 是仿射簇 $X$ 上有限秩局部自由 $\mathcal{O}_X$-模。则 $\Gamma(X,V)$ 是投射 $R(X)$-模。
\end{lemma}

\begin{proof}
首先注意 $\Gamma(X,V)$ 是有限生成 $R(X)$-模。因此由 I.6.1 的类似结论，对任意 $R(X)$-模 $N$ 及任意 $x\in X$ 有
$$\operatorname{Ext}^1_{R(X)}(\Gamma(X,V),N)_x=\operatorname{Ext}^1_{\mathcal{O}_{X,x}}(V_x,N_x).$$
由假设，每个 $V_x$ 是自由 $\mathcal{O}_{X,x}$-模。故对任意 $x$（原文误写为 $x\in R(X)$）有
$\operatorname{Ext}^1_{R(X)}(\Gamma(X,V),N)_x=0$，这意味着
$\operatorname{Ext}^1_{R(X)}(\Gamma(X,V),N)=0$，即 $\Gamma(X,V)$ 是投射 $R(X)$-模。
\end{proof}

\begin{corollary}\label{cor:4.2.4}
设 $X$ 是光滑仿射簇。则 $D_X$ 关于左乘是投射 $R(X)$-模。
\end{corollary}

\begin{proof}
设 $p\in\mathbb{N}$。则有 $\mathcal{O}_X$-模的正合列
$$0\longrightarrow F_{p-1}D_X\longrightarrow F_pD_X\longrightarrow \operatorname{Gr}_pD_X\longrightarrow 0.$$
由于 $X$ 仿射，这给出 $R(X)$-模的正合列
$$0\longrightarrow\Gamma(X,F_{p-1}D_X)\longrightarrow\Gamma(X,F_pD_X)
\longrightarrow\Gamma(X,\operatorname{Gr}_pD_X)\longrightarrow 0.$$
由前述，$\operatorname{Gr}_pD_X$ 是有限秩局部自由 $\mathcal{O}_X$-模。因此由引理 \ref{lem:4.2.3}，$\Gamma(X,\operatorname{Gr}_pD_X)$ 是投射 $R(X)$-模。上述正合列分裂，故
$$\Gamma(X,F_pD_X)=\Gamma(X,F_{p-1}D_X)\oplus\Gamma(X,\operatorname{Gr}_pD_X).$$
由归纳得
$$\Gamma(X,F_pD_X)=\bigoplus_{i=0}^p\Gamma(X,\operatorname{Gr}_iD_X),$$
以及作为 $R(X)$-模
$$D_X=\Gamma(X,\mathcal{D}_X)=\bigoplus_{i=0}^{\infty}\Gamma(X,\operatorname{Gr}_iD_X).$$
因此 $D_X$ 是投射 $R(X)$-模的直和，从而是投射的。
\end{proof}

\begin{lemma}\label{lem:4.2.5}
设 $\varphi:X\to Y$ 是光滑仿射簇之间的态射。则
$D_{X\to Y}=R(X)\otimes_{R(Y)}D_Y$ 是投射 $R(X)$-模。
\end{lemma}

\begin{proof}
由推论 \ref{cor:4.2.4}，$D_Y$ 是投射 $R(Y)$-模，故它是某个自由 $R(Y)$-模的直和因子。于是
$D_{X\to Y}=R(X)\otimes_{R(Y)}D_Y$ 是自由 $R(X)$-模的直和因子，从而是投射 $R(X)$-模。
\end{proof}

\begin{proposition}\label{prop:4.2.6}
设 $X$、$Y$、$Z$ 是光滑仿射簇，$\varphi:X\to Y$、$\psi:Y\to Z$ 是仿射簇的态射。则对 $j\in\mathbb{N}$ 有
$$\operatorname{Tor}^{D_Y}_j(D_{X\to Y},D_{Y\to Z})=0.$$
\end{proposition}

\begin{proof}
设 $M$ 是左 $D_Y$-模，$F^\cdot$ 是它由自由 $D_Y$-模构成的左分解。由推论 \ref{cor:4.2.4}，$D_Y$ 关于左乘是投射 $R(Y)$-模，故也可把 $F^\cdot$ 看作由投射 $R(Y)$-模构成的分解。这意味着
\begin{align*}
\operatorname{Tor}^{R(Y)}_j(R(X),M)&=H^{-j}(R(X)\otimes_{R(Y)}F^\cdot)\\
&=H^{-j}\bigl((R(X)\otimes_{R(Y)}D_Y)\otimes_{D_Y}F^\cdot\bigr)
=\operatorname{Tor}^{D_Y}_j(D_{X\to Y},M).
\end{align*}
由于由引理 \ref{lem:4.2.5}，$D_{Y\to Z}$ 是投射 $R(Y)$-模，故对 $j\in\mathbb{N}$ 有
$\operatorname{Tor}^{R(Y)}_j(R(X),D_{Y\to Z})=0$，结论随之成立。
\end{proof}

下一个结果把 I.12.2 与 I.12.6 推广到光滑仿射簇。

\begin{theorem}\label{thm:4.2.7}
设 $X$、$Y$、$Z$ 是三个光滑仿射簇，$\varphi:X\to Y$、$\psi:Y\to Z$ 是簇的态射。则
\begin{enumerate}[label=(\roman*)]
  \item 对任意左 $D_Z$-模 $M$，存在谱序列，其 $E_2$ 项为
    $E_2^{pq}=L_p\varphi^+(L_q\psi^+(M))$，收敛到 $L_{p+q}(\psi\circ\varphi)^+(M)$；
  \item 对任意右 $D_X$-模 $N$，存在谱序列，其 $E_2$ 项为
    $E_2^{pq}=L_p\psi_+(L_q\varphi_+(N))$，收敛到 $L_{p+q}(\psi\circ\varphi)_+(N)$。
\end{enumerate}
\end{theorem}

\begin{proof}
两个结论都可由 Grothendieck 谱序列得出。
\begin{enumerate}[label=(\roman*)]
  \item 设 $P$ 是投射左 $D_Z$-模，则它是某个自由 $D_Z$-模的直和因子。因此 $\psi^+(P)$ 是某个
    $D_{Y\to Z}^{(I)}$ 的直和因子。由命题 \ref{prop:4.2.6}，这意味着 $\psi^+(P)$ 是 $\varphi^+$-无循环的。
  \item 设 $Q$ 是投射右 $D_X$-模，则它是某个自由 $D_X$-模的直和因子。因此 $\varphi_+(Q)$ 是某个
    $D_{X\to Y}^{(J)}$ 的直和因子。再次应用命题 \ref{prop:4.2.6}，得 $\varphi_+(Q)$ 是 $\psi_+$-无循环的。
\end{enumerate}
\end{proof}

\section{逆像函子}

设 $X$、$Y$ 是两个光滑代数簇，$\varphi:X\to Y$ 是代数簇的态射。定义函子
$$\varphi^+(\mathcal{V})=\mathcal{D}_{X\to Y}\otimes_{\varphi^{-1}(\mathcal{D}_Y)}\varphi^{-1}(\mathcal{V})$$
从左 $\mathcal{D}_Y$-模范畴 $\mathcal{M}_L(\mathcal{D}_Y)$ 到左 $\mathcal{D}_X$-模范畴 $\mathcal{M}_L(\mathcal{D}_X)$。这个函子称为\emph{逆像函子}。由于 $\varphi^{-1}$ 正合而 $\mathcal{D}_{X\to Y}\otimes_{\varphi^{-1}(\mathcal{D}_Y)}-$ 右正合，逆像函子是右正合的。

\begin{remark}\label{rem:4.3.1}
设 $U$ 是 $X$ 的开子集。由注记 \ref{rem:4.1.6} 知 $\mathcal{D}_{U\to X}=\mathcal{D}_U$。故立刻得到
$i^+(\mathcal{V})=i^{-1}(\mathcal{V})=\mathcal{V}|_U$，即 $i$ 的逆像就是限制到 $U$ 的函子
$\mathcal{V}\mapsto\mathcal{V}|_U$。特别地，它是正合函子。
\end{remark}

\begin{lemma}\label{lem:4.3.2}
平坦左 $\mathcal{D}_Y$-模的类对 $\varphi^+$ 是左适合的（left adapted）。
\end{lemma}

\begin{proof}
平坦左 $\mathcal{D}_Y$-模的类对直和封闭，且每个左 $\mathcal{D}_Y$-模都是某个平坦模的商。

设 $F$ 是平坦左 $\mathcal{D}_Y$-模。则对任意 $y\in Y$，$F_y$ 是平坦左 $\mathcal{D}_{Y,y}$-模。因此对任意 $x\in X$，
$\varphi^{-1}(F)_x=F_{\varphi(x)}$ 在 $\varphi^{-1}(\mathcal{D}_Y)_x=\mathcal{D}_{Y,\varphi(x)}$ 上平坦。故
$\varphi^{-1}(F)$ 是平坦左 $\varphi^{-1}(\mathcal{D}_Y)$-模。

因此，若 $F^\cdot$ 是上有界的平坦左 $\mathcal{D}_Y$-模的无循环复形，则 $\varphi^{-1}(F^\cdot)$ 是上有界的平坦左
$\varphi^{-1}(\mathcal{D}_Y)$-模的无循环复形。于是 $\varphi^+(F^\cdot)$ 是上有界无循环复形。

这表明所有平坦左 $\mathcal{D}_Y$-模的类对 $\varphi^+$ 是左适合的。
\end{proof}

因此可定义左导出函子
$$L\varphi^+:D^-(\mathcal{M}_L(\mathcal{D}_Y))\longrightarrow D^-(\mathcal{M}_L(\mathcal{D}_X)).$$

\begin{lemma}\label{lem:4.3.3}
函子 $\varphi^+$ 的左同调维数有限。
\end{lemma}

\begin{proof}
设 $\mathcal{V}$ 是左 $\mathcal{D}_Y$-模，$F^\cdot$ 是它由平坦左 $\mathcal{D}_Y$-模构成的左分解。则由引理 \ref{lem:4.3.2}，
$$L_p\varphi^+(\mathcal{V})=H^p(\varphi^+(F^\cdot))
=H^p\bigl(\mathcal{D}_{X\to Y}\otimes_{\varphi^{-1}(\mathcal{D}_Y)}\varphi^{-1}(F^\cdot)\bigr).$$
设 $x\in X$，则 $F_{\varphi(x)}^\cdot$ 是 $\mathcal{V}_{\varphi(x)}$ 的 $\mathcal{D}_{Y,\varphi(x)}$-平坦分解。故
$$(L_p\varphi^+(\mathcal{V}))_x
=H^p\bigl(\mathcal{D}_{X\to Y}\otimes_{\varphi^{-1}(\mathcal{D}_Y)}\varphi^{-1}(F^\cdot)\bigr)_x
=H^p\bigl(\mathcal{D}_{X\to Y,x}\otimes_{\varphi^{-1}(\mathcal{D}_Y)_x}\varphi^{-1}(F^\cdot)_x\bigr)$$
$$=H^p\bigl(\mathcal{D}_{X\to Y,x}\otimes_{\mathcal{D}_{Y,\varphi(x)}}F_{\varphi(x)}^\cdot\bigr)
=\operatorname{Tor}^{-\,p}_{\mathcal{D}_{Y,\varphi(x)}}(\mathcal{D}_{X\to Y,x},\mathcal{V}_{\varphi(x)}).$$
由于环 $\mathcal{D}_{Y,\varphi(x)}$ 的同调维数 $\le 2\dim Y$，可知当 $p<-2\dim Y$ 时
$(L_p\varphi^+(\mathcal{V}))_x=0$。又因 $x\in X$ 任意，故当 $p<-2\dim Y$ 时 $L_p\varphi^+=0$。
\end{proof}

因此，左导出函子 $L\varphi^+$ 可延拓为无界复形导出范畴之间的导出函子
$$L\varphi^+:D(\mathcal{M}_L(\mathcal{D}_Y))\longrightarrow D(\mathcal{M}_L(\mathcal{D}_X)).$$

\begin{lemma}\label{lem:4.3.4}
设 $\mathcal{V}^\cdot$ 是左 $\mathcal{D}_Y$-模的复形。则
$$L\varphi^+(\mathcal{V}^\cdot)=\mathcal{D}_{X\to Y}\overset{L}{\otimes}_{\varphi^{-1}(\mathcal{D}_Y)}\varphi^{-1}(\mathcal{V}^\cdot).$$
\end{lemma}

\begin{proof}
设 $F$ 是平坦左 $\mathcal{D}_Y$-模。如引理 \ref{lem:4.3.2} 证明中所述，$\varphi^{-1}(F)$ 是平坦左
$\varphi^{-1}(\mathcal{D}_Y)$-模，因而对函子 $\mathcal{D}_{X\to Y}\otimes_{\varphi^{-1}(\mathcal{D}_Y)}-$ 是无循环的。
\end{proof}

设 $\mathcal{V}$ 是左 $\mathcal{D}_Y$-模。则
$$\varphi^+(\mathcal{V})=\mathcal{D}_{X\to Y}\otimes_{\varphi^{-1}(\mathcal{D}_Y)}\varphi^{-1}(\mathcal{V})
=\bigl(\mathcal{O}_X\otimes_{\varphi^{-1}(\mathcal{O}_Y)}\varphi^{-1}(\mathcal{D}_Y)\bigr)
\otimes_{\varphi^{-1}(\mathcal{D}_Y)}\varphi^{-1}(\mathcal{V})$$
$$=\mathcal{O}_X\otimes_{\varphi^{-1}(\mathcal{O}_Y)}
\bigl(\varphi^{-1}(\mathcal{D}_Y)\otimes_{\varphi^{-1}(\mathcal{D}_Y)}\varphi^{-1}(\mathcal{V})\bigr)
=\mathcal{O}_X\otimes_{\varphi^{-1}(\mathcal{O}_Y)}\varphi^{-1}(\mathcal{V})=\varphi^*(\mathcal{V}),$$
即若忘掉 $\mathcal{D}_X$-模结构，则 D-模逆像 $\varphi^+(\mathcal{V})$ 等于 O-模逆像 $\varphi^*(\mathcal{V})$。现证明这对导出函子仍成立。先需要以下结果。

\begin{lemma}\label{lem:4.3.5}
设 $F$ 是平坦左 $\mathcal{D}_Y$-模。则 $F$ 是平坦 $\mathcal{O}_Y$-模。
\end{lemma}

\begin{proof}
设 $y\in Y$，只需证 $F_y$ 是平坦 $\mathcal{O}_{Y,y}$-模。由假设，$F_y$ 是平坦左 $\mathcal{D}_{Y,y}$-模。设 $W$ 是右
$\mathcal{O}_{Y,y}$-模。则
$$W\otimes_{\mathcal{O}_{Y,y}}F_y=W\otimes_{\mathcal{O}_{Y,y}}(\mathcal{D}_{Y,y}\otimes_{\mathcal{D}_{Y,y}}F_y)
=(W\otimes_{\mathcal{O}_{Y,y}}\mathcal{D}_{Y,y})\otimes_{\mathcal{D}_{Y,y}}F_y.$$
由于 $\mathcal{D}_{Y,y}$ 关于左乘是自由 $\mathcal{O}_{Y,y}$-模，故 $W\mapsto W\otimes_{\mathcal{O}_{Y,y}}\mathcal{D}_{Y,y}$ 是正合函子。于是
$W\mapsto(W\otimes_{\mathcal{O}_{Y,y}}\mathcal{D}_{Y,y})\otimes_{\mathcal{D}_{Y,y}}F_y$ 也是正合的。这立即使
$W\mapsto W\otimes_{\mathcal{O}_{Y,y}}F_y$ 正合，即 $F_y$ 是 $\mathcal{O}_{Y,y}$-平坦的。
\end{proof}

这导致以下结果。

\begin{theorem}\label{thm:4.3.6}
若竖直箭头表示遗忘函子，则如下函子图
$$\begin{tikzcd}
D(\mathcal{M}_L(\mathcal{D}_Y)) \arrow[r,"L\varphi^+"] \arrow[d,"\operatorname{For}"'] &
D(\mathcal{M}_L(\mathcal{D}_X)) \arrow[d,"\operatorname{For}"]\\
D(\mathcal{M}(\mathcal{O}_Y)) \arrow[r,"L\varphi^*"'] & D(\mathcal{M}(\mathcal{O}_X))
\end{tikzcd}$$
在同构意义下交换。
\end{theorem}

\begin{proof}
这可由引理 \ref{lem:4.3.4} 以及关于函子复合的导出函子定理，分别应用于遗忘函子与 $\varphi^+$ 的复合、$\varphi^*$ 与遗忘函子的复合而得出。
\end{proof}

这个结果结合前述结论，有以下直接推论。

\begin{theorem}\label{thm:4.3.7}
设 $\mathcal{V}^\cdot$ 是左 $\mathcal{D}_Y$-模的复形，且对所有 $p\in\mathbb{Z}$，$H^p(\mathcal{V}^\cdot)$ 是拟凝聚的。则对所有 $p\in\mathbb{Z}$，上同调模 $H^p(L\varphi^+(\mathcal{V}^\cdot))$ 是拟凝聚左 $\mathcal{D}_X$-模。
\end{theorem}

现设 $X$、$Y$ 是光滑仿射簇。一般地，对任意左 $\mathcal{D}_Y$-模 $\mathcal{V}$，有自然态射
$$\mathcal{D}_{X\to Y}\otimes_{\mathcal{D}_Y}\Gamma(Y,\mathcal{V})\longrightarrow\Gamma(X,\varphi^+(\mathcal{V})).$$
下一个结果蕴含：对拟凝聚的 $\mathcal{V}$，这个态射是同构。

\begin{proposition}\label{prop:4.3.8}
设 $X$、$Y$ 是光滑仿射簇，$\varphi:X\to Y$ 是代数簇的态射。则对任意拟凝聚左 $\mathcal{D}_Y$-模 $\mathcal{V}$，有
$$\Gamma(X,L_p\varphi^+(\mathcal{V}))=\operatorname{Tor}^{-\,p}_{\mathcal{D}_Y}(D_{X\to Y},\Gamma(Y,\mathcal{V}))$$
对 $p\in\mathbb{Z}$ 成立。
\end{proposition}

\begin{proof}
首先注意若 $\mathcal{V}=\mathcal{D}_Y$，则 $\varphi^+(\mathcal{D}_Y)=\mathcal{D}_{X\to Y}$，上述态射显然是同构。现考虑任意拟凝聚左 $\mathcal{D}_Y$-模 $\mathcal{V}$。设 $F^\cdot$ 是 $\Gamma(Y,\mathcal{V})$ 的自由左分解。则它的局部化 $\mathcal{F}^\cdot=\Delta(F^\cdot)$ 是 $\mathcal{V}$ 的自由分解，即 $\mathcal{V}$ 的 $\mathcal{D}_Y$-平坦分解。因此
$$L_p\varphi^+(\mathcal{V})=H^p(L\varphi^+(\mathcal{D}(\mathcal{V})))=H^p(\varphi^+(\mathcal{F}^\cdot))$$
对 $p\in\mathbb{Z}$ 成立。由于 $\varphi^+(\mathcal{F}^\cdot)$ 是拟凝聚 $\mathcal{D}_X$-模的复形，
$$\Gamma(X,L_p\varphi^+(\mathcal{V}))=\Gamma(X,H^p(\varphi^+(\mathcal{F}^\cdot)))=H^p(\Gamma(X,\varphi^+(\mathcal{F}^\cdot))).$$
最后，由于 $\Gamma(X,-)$ 与 $\varphi^+$ 都与直和可交换，由证明的第一部分得
$$\Gamma(X,L_p\varphi^+(\mathcal{V}))=H^p(D_{X\to Y}\otimes_{\mathcal{D}_Y}F^\cdot)
=\operatorname{Tor}^{-\,p}_{\mathcal{D}_Y}(D_{X\to Y},\Gamma(Y,\mathcal{V}))$$
对 $p\in\mathbb{Z}$ 成立。
\end{proof}

因此，在光滑仿射簇的态射与拟凝聚 D-模的情形下，我们恢复了 §2 中的旧定义。

设 $X$、$Y$、$Z$ 是三个光滑代数簇，$\varphi:X\to Y$、$\psi:Y\to Z$ 是两个簇的态射。则对任意左 $\mathcal{D}_Z$-模 $\mathcal{V}$，有
$$(\varphi^+\circ\psi^+)(\mathcal{V})=\mathcal{D}_{X\to Y}\otimes_{\varphi^{-1}(\mathcal{D}_Y)}\varphi^{-1}(\psi^+(\mathcal{V}))$$
$$=\mathcal{D}_{X\to Y}\otimes_{\varphi^{-1}(\mathcal{D}_Y)}\varphi^{-1}
(\mathcal{D}_{Y\to Z}\otimes_{\psi^{-1}(\mathcal{D}_Z)}\psi^{-1}(\mathcal{V}))$$
$$=\mathcal{D}_{X\to Y}\otimes_{\varphi^{-1}(\mathcal{D}_Y)}
\bigl(\varphi^{-1}(\mathcal{D}_{Y\to Z})\otimes_{\varphi^{-1}(\psi^{-1}(\mathcal{D}_Z))}\varphi^{-1}(\psi^{-1}(\mathcal{V}))\bigr)$$
$$=\bigl(\mathcal{D}_{X\to Y}\otimes_{\varphi^{-1}(\mathcal{D}_Y)}\varphi^{-1}(\mathcal{D}_{Y\to Z})\bigr)
\otimes_{(\psi\circ\varphi)^{-1}(\mathcal{D}_Z)}(\psi\circ\varphi)^{-1}(\mathcal{V}).$$
故由前述，
$$(\varphi^+\circ\psi^+)(\mathcal{V})=\mathcal{D}_{X\to Z}
\otimes_{(\psi\circ\varphi)^{-1}(\mathcal{D}_Z)}(\psi\circ\varphi)^{-1}(\mathcal{V})=(\psi\circ\varphi)^+(\mathcal{V}).$$
因此我们证明了以下结果。

\begin{lemma}\label{lem:4.3.9}
$\varphi^+\circ\psi^+=(\psi\circ\varphi)^+$。
\end{lemma}

下一个结果把它推广到导出范畴。

\begin{theorem}\label{thm:4.3.10}
正合函子 $L\varphi^+\circ L\psi^+$ 与 $L(\psi\circ\varphi)^+$（从 $D(\mathcal{M}_L(\mathcal{D}_Z))$ 到
$D(\mathcal{M}_L(\mathcal{D}_X))$）是同构的。
\end{theorem}

\begin{proof}
由上述讨论及注记 \ref{rem:4.3.1}，只需验证对任意平坦左 $\mathcal{D}_Z$-模 $F$，$\psi^+(F)$ 是 $\varphi^+$-无循环的。由引理 \ref{lem:4.3.5}，只需证 $\psi^*(F)$ 是 $\varphi^*$-无循环的。这由引理 \ref{lem:4.3.4} 及前述讨论得出。
\end{proof}

显然有
$$\varphi^+(\mathcal{D}_Y)=\mathcal{D}_{X\to Y}\otimes_{\varphi^{-1}(\mathcal{D}_Y)}\varphi^{-1}(\mathcal{D}_Y)=\mathcal{D}_{X\to Y}.$$
由于 $\mathcal{D}_Y$ 是平坦左 $\mathcal{D}_Y$-模，由注记 \ref{rem:4.3.1} 得
$$L\varphi^+(\mathcal{D}(\mathcal{D}_Y))=\mathcal{D}(\mathcal{D}_{X\to Y}).$$
故由引理 \ref{lem:4.3.9} 得
$$\mathcal{D}(\mathcal{D}_{X\to Z})=L(\psi\circ\varphi)^+(\mathcal{D}(\mathcal{D}_Z))
=L\varphi^+(L\psi^+(\mathcal{D}(\mathcal{D}_Z)))$$
$$=L\varphi^+(\mathcal{D}(\mathcal{D}_{Y\to Z}))
=\mathcal{D}_{X\to Y}\overset{L}{\otimes}_{\varphi^{-1}(\mathcal{D}_Y)}\mathcal{D}(\varphi^{-1}(\mathcal{D}_{Y\to Z})).$$

\begin{corollary}\label{cor:4.3.11}
有
$$\mathcal{D}(\mathcal{D}_{X\to Z})=\mathcal{D}(\mathcal{D}_{X\to Y})
\overset{L}{\otimes}_{\varphi^{-1}(\mathcal{D}_Y)}\mathcal{D}(\varphi^{-1}(\mathcal{D}_{Y\to Z})).$$
\end{corollary}

\section{投影公式}

设 $X$、$Y$ 是两个拓扑空间，$\varphi:X\to Y$ 是连续映射。设 $R$ 是 $Y$ 上的环层，则 $\varphi^{-1}(R)$ 是 $X$ 上的环层。

设 $A$ 是 $X$ 上右 $\varphi^{-1}(R)$-模，$B$ 是 $Y$ 上左 $R$-模。则可考虑 $Y$ 上的阿贝尔群层
$\varphi_\bullet(A\otimes_{\varphi^{-1}(R)}\varphi^{-1}(B))$ 与 $\varphi_\bullet(A)\otimes_R B$。第一个由
$$V\longmapsto\varphi_\bullet(A\otimes_{\varphi^{-1}(R)}\varphi^{-1}(B))(V)
=(A\otimes_{\varphi^{-1}(R)}\varphi^{-1}(B))(\varphi^{-1}(V))$$
给出，第二个是预层
$$V\longmapsto\varphi_\bullet(A)(V)\otimes_{R(V)}B(V)=A(\varphi^{-1}(V))\otimes_{R(V)}B(V)$$
相伴的层。

由于 $A\otimes_{\varphi^{-1}(R)}\varphi^{-1}(B)$ 是 $X$ 上预层
$$U\longmapsto A(U)\otimes_{\varphi^{-1}(R)(U)}\varphi^{-1}(B)(U)$$
相伴的层，对任意开集 $V\subset Y$，有从
$$A(\varphi^{-1}(V))\otimes_{R(V)}B(V)$$
到
$$A(\varphi^{-1}(V))\otimes_{\varphi^{-1}(R)(\varphi^{-1}(V))}\varphi^{-1}(B)(\varphi^{-1}(V))$$
再到群
$$(A\otimes_{\varphi^{-1}(R)}\varphi^{-1}(B))(\varphi^{-1}(V))$$
的自然态射。因此有自然态射
$$\varphi_\bullet(A)\otimes_R B\longrightarrow\varphi_\bullet(A\otimes_{\varphi^{-1}(R)}\varphi^{-1}(B)).$$
这显然是双函子的态射。因此它诱导有上界复形同伦范畴之间相应双函子的态射：若 $A^\cdot$ 是上有界右 $\varphi^{-1}(R)$-模的复形，$B^\cdot$ 是上有界左 $R$-模的复形，则有典范态射
$$\varphi_\bullet(A^\cdot)\otimes_R B^\cdot\longrightarrow
\varphi_\bullet(A^\cdot\otimes_{\varphi^{-1}(R)}\varphi^{-1}(B^\cdot)).$$
现设 $\varphi_\bullet$ 具有有限右同调维数。则有典范态射
$$\varphi_\bullet(A^\cdot\otimes_{\varphi^{-1}(R)}\varphi^{-1}(B^\cdot))
\longrightarrow R\varphi_\bullet(A^\cdot\otimes_{\varphi^{-1}(R)}\varphi^{-1}(B^\cdot)).$$
设 $A^\cdot$ 是 $\varphi_\bullet$-无循环的且 $B^\cdot$ 是 $R$-平坦的。则由于 $\varphi^{-1}(B^\cdot)$ 是 $\varphi^{-1}(R)$-平坦的，有典范态射
$$\varphi_\bullet(A^\cdot)\otimes_R B^\cdot\longrightarrow
R\varphi_\bullet(A^\cdot\otimes_{\varphi^{-1}(R)}\varphi^{-1}(B^\cdot))
=R\varphi_\bullet(A^\cdot\overset{L}{\otimes}_{\varphi^{-1}(R)}\varphi^{-1}(B^\cdot)).$$
最后，由于任意上有界右 $\varphi^{-1}(R)$-模复形拟同构于某个上有界 $\varphi_\bullet$-无循环复形，任意上有界左 $R$-模复形拟同构于某个上有界 $R$-平坦复形，得到自然态射
$$R\varphi_\bullet(A^\cdot)\overset{L}{\otimes}_R B^\cdot\longrightarrow
R\varphi_\bullet(A^\cdot\overset{L}{\otimes}_{\varphi^{-1}(R)}\varphi^{-1}(B^\cdot)).$$
因此在导出范畴之间得到了相应双函子的态射。

在某些条件下，这个自然态射是同构。相应的陈述通常称为\emph{投影公式}。例如，设 $X$、$Y$ 是代数簇，$\varphi:X\to Y$ 是代数簇的态射。此时 $\varphi_\bullet$ 具有有限右同调维数。

\begin{proposition}\label{prop:4.4.1}
设 $X$、$Y$ 是两个代数簇，$\varphi:X\to Y$ 是代数簇的态射。则
$$R\varphi_\bullet(\mathcal{V}^\cdot)\overset{L}{\otimes}_{\mathcal{O}_Y}\mathcal{W}^\cdot
=R\varphi_\bullet(\mathcal{V}^\cdot\overset{L}{\otimes}_{\varphi^{-1}(\mathcal{O}_Y)}\varphi^{-1}(\mathcal{W}^\cdot))$$
对任意 $D^-(\mathcal{M}(\varphi^{-1}(\mathcal{O}_Y)))$ 中的 $\mathcal{V}^\cdot$ 与
$D^-(\mathcal{M}_{qc}(\mathcal{O}_Y))$ 中的 $\mathcal{W}^\cdot$ 成立。
\end{proposition}

\begin{proof}
只需证明典范态射在上同调群上诱导同构。由于 $\varphi_\bullet$ 的右同调维数有限，通过截断点可设 $\mathcal{W}^\cdot$ 是有界复形。此外，由于所有拟凝聚 $\mathcal{O}_Y$-模构成有界导出范畴的一个生成类，可设 $\mathcal{W}^\cdot=\mathcal{D}(\mathcal{W})$，其中 $\mathcal{W}$ 是拟凝聚 $\mathcal{O}_Y$-模。

由于该陈述关于 $Y$ 是局部的，还可设 $Y$ 是仿射簇。此时可用自由分解 $F^\cdot$ 代替 $\mathcal{W}$。由于 $\varphi_\bullet$ 与直和可交换，显然此时自然态射是同构。
\end{proof}

类似地可证明以下陈述。

\begin{proposition}\label{prop:4.4.2}
设 $X$、$Y$ 是两个代数簇，$\varphi:X\to Y$ 是代数簇的态射。则
$$R\varphi_\bullet(\mathcal{V}^\cdot)\overset{L}{\otimes}_{\mathcal{D}_Y}\mathcal{W}^\cdot
=R\varphi_\bullet(\mathcal{V}^\cdot\overset{L}{\otimes}_{\varphi^{-1}(\mathcal{D}_Y)}\varphi^{-1}(\mathcal{W}^\cdot))$$
对任意 $D^-(\mathcal{M}_R(\varphi^{-1}(\mathcal{D}_Y)))$ 中的 $\mathcal{V}^\cdot$ 与
$D^-(\mathcal{M}^L_{qc}(\mathcal{D}_Y))$ 中的 $\mathcal{W}^\cdot$ 成立。
\end{proposition}

设 $Y$ 是光滑簇。则 $\mathcal{O}_{Y,y}$ 的同调维数 $\le\dim Y$，$\mathcal{D}_{Y,y}$ 的同调维数 $\le 2\dim Y$。因此用标准的截断点论证，可建立前两个结果的下述变体。

\begin{proposition}\label{prop:4.4.3}
设 $X$、$Y$ 是两个光滑代数簇，$\varphi:X\to Y$ 是代数簇的态射。则
$$R\varphi_\bullet(\mathcal{V}^\cdot)\overset{L}{\otimes}_{\mathcal{O}_Y}\mathcal{W}^\cdot
=R\varphi_\bullet(\mathcal{V}^\cdot\overset{L}{\otimes}_{\varphi^{-1}(\mathcal{O}_Y)}\varphi^{-1}(\mathcal{W}^\cdot))$$
对任意 $D(\mathcal{M}(\varphi^{-1}(\mathcal{O}_Y)))$ 中的 $\mathcal{V}^\cdot$ 与
$D^b(\mathcal{M}_{qc}(\mathcal{O}_Y))$ 中的 $\mathcal{W}^\cdot$ 成立。
\end{proposition}

\begin{proposition}\label{prop:4.4.4}
设 $X$、$Y$ 是两个光滑代数簇，$\varphi:X\to Y$ 是代数簇的态射。则
$$R\varphi_\bullet(\mathcal{V}^\cdot)\overset{L}{\otimes}_{\mathcal{D}_Y}\mathcal{W}^\cdot
=R\varphi_\bullet(\mathcal{V}^\cdot\overset{L}{\otimes}_{\varphi^{-1}(\mathcal{D}_Y)}\varphi^{-1}(\mathcal{W}^\cdot))$$
对任意 $D(\mathcal{M}_R(\varphi^{-1}(\mathcal{D}_Y)))$ 中的 $\mathcal{V}^\cdot$ 与
$D^b(\mathcal{M}^L_{qc}(\mathcal{D}_Y))$ 中的 $\mathcal{W}^\cdot$ 成立。
\end{proposition}

特别地，我们需要以下形式的投影公式。设 $\varphi:X\to Y$ 与 $\psi:Y\to Z$ 是光滑代数簇的态射。

\begin{lemma}\label{lem:4.4.5}
设 $\mathcal{V}^\cdot$ 是右 $\varphi^{-1}(\mathcal{D}_Y)$-模的复形。则在
$D(\mathcal{M}_R(\psi^{-1}(\mathcal{D}_Z)))$ 中有右 $\psi^{-1}(\mathcal{D}_Z)$-模复形的自然同构
$$R\varphi_\bullet(\mathcal{V}^\cdot)\overset{L}{\otimes}_{\mathcal{D}_Y}\mathcal{D}_{Y\to Z}
=R\varphi_\bullet(\mathcal{V}^\cdot\overset{L}{\otimes}_{\varphi^{-1}(\mathcal{D}_Y)}
\varphi^{-1}(\mathcal{D}_{Y\to Z})).$$
\end{lemma}

\begin{proof}
可把双模 $\mathcal{D}_{Y\to Z}$ 看作环层 $\mathcal{D}_Y\otimes_k\psi^{-1}(\mathcal{D}_Z^{opp})$ 上的模层，其中
$\mathcal{D}_Z^{opp}$ 是 $\mathcal{D}_Z$ 的对立环层。显然，平坦的
$\mathcal{D}_Y\otimes_k\psi^{-1}(\mathcal{D}_Z^{opp})$-模作为 $\mathcal{D}_Y$-模是平坦的。因此从
$D(\mathcal{M}_R(\varphi^{-1}(\mathcal{D}_Y)))$ 到 $D(\mathcal{M}_R(\psi^{-1}(\mathcal{D}_Z)))$ 的函子典范态射
$$R\varphi_\bullet(\mathcal{V}^\cdot)\overset{L}{\otimes}_{\mathcal{D}_Y}\mathcal{D}_{Y\to Z}
\longrightarrow R\varphi_\bullet(\mathcal{V}^\cdot\overset{L}{\otimes}_{\varphi^{-1}(\mathcal{D}_Y)}
\varphi^{-1}(\mathcal{D}_{Y\to Z}))$$
若忘掉 $\psi^{-1}(\mathcal{D}_Z^{opp})$-模作用，即诱导出命题 \ref{prop:4.4.4} 的典范态射。另一方面，由命题 \ref{prop:4.4.4}，这个态射是同构。
\end{proof}

\section{直像函子}

设 $X$、$Y$ 是两个光滑代数簇，$\varphi:X\to Y$ 是代数簇的态射。由于对 $x\in X$，$\mathcal{D}_{X,x}$ 的同调维数 $\le 2\dim X$，函子 $U\mapsto U\otimes_{\mathcal{D}_X}\mathcal{D}_{X\to Y}$ 具有有限左同调维数。因此有函子
$\mathcal{V}^\cdot\mapsto\mathcal{V}^\cdot\overset{L}{\otimes}_{\mathcal{D}_X}\mathcal{D}_{X\to Y}$，介于
$D(\mathcal{M}_R(\mathcal{D}_X))$ 与 $D(\mathcal{M}_R(\varphi^{-1}(\mathcal{D}_Y)))$ 之间。另一方面，$\varphi_\bullet$ 也具有有限同调维数，故有函子
$\mathcal{V}^\cdot\mapsto R\varphi_\bullet(\mathcal{V}^\cdot\overset{L}{\otimes}_{\mathcal{D}_X}\mathcal{D}_{X\to Y})$，从
$D(\mathcal{M}_R(\mathcal{D}_X))$ 到 $D(\mathcal{M}_R(\mathcal{D}_Y))$。我们称这个函子为直像函子
$\varphi_+:D(\mathcal{M}_R(\mathcal{D}_X))\to D(\mathcal{M}_R(\mathcal{D}_Y))$，特别地
$$\varphi_+(\mathcal{V}^\cdot)=R\varphi_\bullet(\mathcal{V}^\cdot\overset{L}{\otimes}_{\mathcal{D}_X}\mathcal{D}_{X\to Y})$$
对任意 $D(\mathcal{M}_R(\mathcal{D}_X))$ 中的 $\mathcal{V}^\cdot$ 成立。

\begin{lemma}\label{lem:4.5.1}
正合函子 $\varphi_+:D(\mathcal{M}_R(\mathcal{D}_X))\to D(\mathcal{M}_R(\mathcal{D}_Y))$ 具有有限振幅（amplitude）。
\end{lemma}

\begin{proof}
由于 $U\mapsto U\otimes_{\mathcal{D}_X}\mathcal{D}_{X\to Y}$ 具有有限左同调维数，函子
$\mathcal{V}^\cdot\mapsto\mathcal{V}^\cdot\overset{L}{\otimes}_{\mathcal{D}_X}\mathcal{D}_{X\to Y}$ 具有有限振幅。类似地，$R\varphi_\bullet$ 具有有限振幅。这意味着它们的复合具有有限振幅。
\end{proof}

\begin{remark}\label{rem:4.5.2}
设 $U$ 是 $X$ 中开集，$i:U\to X$ 是自然包含。由注记 \ref{rem:4.1.6} 有 $\mathcal{D}_{U\to X}=\mathcal{D}_U$。故此时对任意右 $\mathcal{D}_U$-模复形 $\mathcal{V}^\cdot$ 有
$$i_+(\mathcal{V}^\cdot)=Ri_\bullet(\mathcal{V}^\cdot).$$
特别地，在开包含的情形下，直像函子等于层直像函子 $i_\bullet$ 的右导出函子。
\end{remark}

设 $Z$ 是另一个光滑簇，$\psi:Y\to Z$ 是代数簇的态射。

\begin{theorem}\label{thm:4.5.3}
正合函子 $\psi_+\circ\varphi_+$ 与 $(\psi\circ\varphi)_+$（从 $D(\mathcal{M}_R(\mathcal{D}_X))$ 到
$D(\mathcal{M}_R(\mathcal{D}_Z))$）是同构的。
\end{theorem}

\begin{proof}
设 $\mathcal{V}^\cdot$ 是 $D(\mathcal{M}_R(\mathcal{D}_X))$ 中的复形。则
$$(\psi_+\circ\varphi_+)(\mathcal{V}^\cdot)=\psi_+(\varphi_+(\mathcal{V}^\cdot))
=R\psi_\bullet(\varphi_+(\mathcal{V}^\cdot)\overset{L}{\otimes}_{\mathcal{D}_Y}\mathcal{D}_{Y\to Z})$$
$$=R\psi_\bullet\bigl(R\varphi_\bullet(\mathcal{V}^\cdot\overset{L}{\otimes}_{\mathcal{D}_X}\mathcal{D}_{X\to Y})
\overset{L}{\otimes}_{\mathcal{D}_Y}\mathcal{D}_{Y\to Z}\bigr).$$
由投影公式，
$$R\varphi_\bullet(\mathcal{V}^\cdot\overset{L}{\otimes}_{\mathcal{D}_X}\mathcal{D}_{X\to Y})
\overset{L}{\otimes}_{\mathcal{D}_Y}\mathcal{D}_{Y\to Z}
=R\varphi_\bullet\bigl((\mathcal{V}^\cdot\overset{L}{\otimes}_{\mathcal{D}_X}\mathcal{D}_{X\to Y})
\overset{L}{\otimes}_{\varphi^{-1}(\mathcal{D}_Y)}\varphi^{-1}(\mathcal{D}_{Y\to Z})\bigr).$$
故由前述，
$$(\psi_+\circ\varphi_+)(\mathcal{V}^\cdot)
=R\psi_\bullet\bigl(R\varphi_\bullet(\mathcal{V}^\cdot\overset{L}{\otimes}_{\mathcal{D}_X}
(\mathcal{D}_{X\to Y}\overset{L}{\otimes}_{\varphi^{-1}(\mathcal{D}_Y)}\varphi^{-1}(\mathcal{D}_{Y\to Z})))\bigr)$$
$$=R(\psi\circ\varphi)_\bullet(\mathcal{V}^\cdot\overset{L}{\otimes}_{\mathcal{D}_X}\mathcal{D}_{X\to Z})
=(\psi\circ\varphi)_+(\mathcal{V}^\cdot).$$
\end{proof}

设 $\varphi:X\to Y$ 是光滑代数簇的态射。设 $V$ 是 $Y$ 中开集，$U=\varphi^{-1}(V)$。设 $\mathcal{V}^\cdot$ 是右 $\mathcal{D}_X$-模的复形。则
$$\varphi_+(\mathcal{V}^\cdot)|_V=R\varphi_\bullet(\mathcal{V}^\cdot\overset{L}{\otimes}_{\mathcal{D}_X}\mathcal{D}_{X\to Y})|_V
=R(\varphi|_U)_\bullet\bigl((\mathcal{V}^\cdot\overset{L}{\otimes}_{\mathcal{D}_X}\mathcal{D}_{X\to Y})|_U\bigr)$$
$$=R(\varphi|_U)_\bullet(\mathcal{V}^\cdot|_U\overset{L}{\otimes}_{\mathcal{D}_U}\mathcal{D}_{U\to V})
=(\varphi|_U)_+(\mathcal{V}^\cdot|_U),$$
即直像函子关于目标簇是局部的。

现设 $X$、$Y$ 是光滑仿射簇。则对任意拟凝聚右 $\mathcal{D}_X$-模 $\mathcal{V}$，设 $F^\cdot$ 是 $\Gamma(X,\mathcal{V})$ 的自由分解。则 $\mathcal{F}^\cdot=\Delta(F^\cdot)$ 是 $\mathcal{V}$ 的自由分解。因此
$$\varphi_+(\mathcal{D}(\mathcal{V}))=R\varphi_\bullet(\mathcal{D}(\mathcal{V})\overset{L}{\otimes}_{\mathcal{D}_X}\mathcal{D}_{X\to Y})
=R\varphi_\bullet(F^\cdot\otimes_{\mathcal{D}_X}\mathcal{D}_{X\to Y}).$$
显然对任意 $p\in\mathbb{Z}$，$F^p\otimes_{\mathcal{D}_X}\mathcal{D}_{X\to Y}$ 是 $\mathcal{D}_{X\to Y}$ 若干拷贝的直和，故它是拟凝聚左 $\mathcal{D}_X$-模。特别地，由前述它对 $\varphi_\bullet$ 是无循环的。这意味着
$$\varphi_+(\mathcal{D}(\mathcal{V}))=\varphi_\bullet(F^\cdot\otimes_{\mathcal{D}_X}\mathcal{D}_{X\to Y}).$$

\begin{lemma}\label{lem:4.5.4}
若 $\varphi:X\to Y$ 是光滑仿射簇的态射，则
\begin{enumerate}[label=(\roman*)]
  \item $H^p(\varphi_+(\mathcal{D}(\mathcal{D}_X)))=0$ 对所有 $p\ne 0$；
  \item $H^0(\varphi_+(\mathcal{D}(\mathcal{D}_X)))=\varphi_\bullet(\mathcal{D}_{X\to Y})$ 是拟凝聚右 $\mathcal{D}_Y$-模。
\end{enumerate}
\end{lemma}

\begin{proof}
显然
$$\varphi_+(\mathcal{D}(\mathcal{D}_X))=\varphi_\bullet(\mathcal{D}(\mathcal{D}_X)\otimes_{\mathcal{D}_X}\mathcal{D}_{X\to Y})
=\mathcal{D}(\varphi_\bullet(\mathcal{D}_{X\to Y})).$$
这蕴含 (i)。只需检验 $H^0(\varphi_+(\mathcal{D}(\mathcal{D}_X)))=\varphi_\bullet(\mathcal{D}_{X\to Y})$ 是拟凝聚的。首先注意
$$\Gamma(Y,\varphi_\bullet(\mathcal{D}_{X\to Y}))=\Gamma(X,\mathcal{D}_{X\to Y})=R(X)\otimes_{R(Y)}D_Y,$$
因为 $\mathcal{D}_{X\to Y}$ 是拟凝聚 $\mathcal{D}_X$-模。设 $g\in R(Y)$、$f=g\circ\varphi$。则 $\varphi^{-1}(Y_g)=X_f$，且由于 $\varphi_+$ 是局部的，
$$H^0(\varphi_+(\mathcal{D}(\mathcal{D}_X)))|_{Y_g}=H^0(\varphi_+(\mathcal{D}(\mathcal{D}_{X_f}))),$$
以及
$$\varphi_\bullet(\mathcal{D}_{X\to Y})(Y_g)=H^0(\varphi_+(\mathcal{D}(\mathcal{D}_X)))(Y_g)
=H^0(\varphi_+(\mathcal{D}(\mathcal{D}_{X_f})))(Y_g)$$
$$=\Gamma(Y_g,(\varphi|_{X_f})_\bullet(\mathcal{D}_{X_f\to Y_g}))
=R(X_f)\otimes_{R(Y_g)}D_{Y_g}=R(X)_f\otimes_{R(Y_g)}D_{Y_g}$$
$$=\bigl(R(X)\otimes_{R(Y)}R(Y_g)\bigr)\otimes_{R(Y_g)}D_{Y_g}
=R(X)\otimes_{R(Y)}D_{Y_g}=\bigl(R(X)\otimes_{R(Y)}D_Y\bigr)_g=(\mathcal{D}_{X\to Y})_g,$$
其中局部化是关于第二个因子中的右乘。这说明 $\varphi_\bullet(\mathcal{D}_{X\to Y})$ 是 $\mathcal{D}_{X\to Y}$ 作为右 $\mathcal{D}_Y$-模的局部化，故它是拟凝聚的。
\end{proof}

因此，$\varphi_\bullet(F^\cdot\otimes_{\mathcal{D}_X}\mathcal{D}_{X\to Y})$ 中的右 $\mathcal{D}_Y$-模是拟凝聚的。这意味着该复形的上同调群也是拟凝聚右 $\mathcal{D}_Y$-模。于是 $H^p(\varphi_+(\mathcal{D}(\mathcal{V})))$ 是拟凝聚右 $\mathcal{D}_Y$-模。此外，
$$\Gamma(Y,H^p(\varphi_+(\mathcal{D}(\mathcal{V}))))=H^p(R\Gamma(\varphi_+(\mathcal{D}(\mathcal{V}))))$$
$$=H^p(R\Gamma(\varphi_\bullet(F^\cdot\otimes_{\mathcal{D}_X}\mathcal{D}_{X\to Y})))
=H^p(R\Gamma(F^\cdot\otimes_{\mathcal{D}_X}\mathcal{D}_{X\to Y}))$$
$$=H^p(F^\cdot\otimes_{\mathcal{D}_X}\mathcal{D}_{X\to Y})
=H^p(\Gamma(X,\mathcal{V})\otimes_{\mathcal{D}_X}\mathcal{D}_{X\to Y})
=\operatorname{Tor}^{-\,p}_{\mathcal{D}_X}(\Gamma(X,\mathcal{V}),\mathcal{D}_{X\to Y}).$$
因此我们证明了以下结果。

\begin{proposition}\label{prop:4.5.5}
设 $\varphi:X\to Y$ 是光滑仿射簇的态射。设 $\mathcal{V}$ 是拟凝聚右 $\mathcal{D}_X$-模。则：
\begin{enumerate}[label=(\roman*)]
  \item 对 $p\in\mathbb{Z}$，$H^p(\varphi_+(\mathcal{D}(\mathcal{V})))$ 是拟凝聚右 $\mathcal{D}_Y$-模；
  \item 对 $p>0$，$H^p(\varphi_+(\mathcal{D}(\mathcal{V})))=0$；
  \item 对 $p<0$，有
    $$\Gamma(Y,H^p(\varphi_+(\mathcal{D}(\mathcal{V}))))
    =\operatorname{Tor}^{-\,p}_{\mathcal{D}_X}(\Gamma(X,\mathcal{V}),\mathcal{D}_{X\to Y}).$$
\end{enumerate}
\end{proposition}

因此，在光滑仿射簇上拟凝聚右 D-模的情形下，我们的定义与前述旧定义一致。

设 $\varphi:X\to Y$ 是仿射态射。于是对 $Y$ 中仿射开集 $V$，集合 $U=\varphi^{-1}(V)$ 也是 $X$ 中仿射开子集。设 $\mathcal{V}$ 是拟凝聚右 $\mathcal{D}_X$-模。则
$$\varphi_+(\mathcal{D}(\mathcal{V}))|_V=(\varphi|_U)_+(\mathcal{D}(\mathcal{V}|_U)).$$
因此前述结果有以下推论。

\begin{corollary}\label{cor:4.5.6}
设 $\varphi:X\to Y$ 是光滑簇之间的仿射态射。设 $\mathcal{V}$ 是拟凝聚右 $\mathcal{D}_X$-模。则：
\begin{enumerate}[label=(\roman*)]
  \item 对 $p\in\mathbb{Z}$，$H^p(\varphi_+(\mathcal{D}(\mathcal{V})))$ 是拟凝聚右 $\mathcal{D}_Y$-模；
  \item 对 $p>0$，$H^p(\varphi_+(\mathcal{D}(\mathcal{V})))=0$。
\end{enumerate}
\end{corollary}

第一个陈述在一般情形下也成立。先考虑上述结果的一个特殊情形。

\begin{corollary}\label{cor:4.5.7}
设 $U$ 是 $X$ 中仿射开集，$i:U\to X$ 是自然包含。设 $\mathcal{V}$ 是拟凝聚右 $\mathcal{D}_U$-模。则对任意拟凝聚右 $\mathcal{D}_U$-模 $\mathcal{V}$ 与 $p\ne 0$ 有 $H^p(i_+(\mathcal{D}(\mathcal{V})))=0$。此外，函子 $H^0\circ i_+\circ\mathcal{D}=i_\bullet$ 从拟凝聚右 $\mathcal{D}_U$-模范畴到右 $\mathcal{D}_X$-模范畴是正合的。
\end{corollary}

现在可以证明上述结果的如下推广。

\begin{theorem}\label{thm:4.5.8}
设 $\varphi:X\to Y$ 是光滑代数簇的态射。设 $\mathcal{V}^\cdot$ 是右 $\mathcal{D}_X$-模的复形，且对所有 $p\in\mathbb{Z}$，$H^p(\mathcal{V}^\cdot)$ 是拟凝聚右 $\mathcal{D}_X$-模。则对所有 $p\in\mathbb{Z}$，$H^p(\varphi_+(\mathcal{V}^\cdot))$ 是拟凝聚右 $\mathcal{D}_Y$-模。
\end{theorem}

\begin{proof}
设 $U$ 是 $X$ 中仿射开子集，$i:U\to X$ 是自然浸入。设 $\mathcal{W}$ 是拟凝聚右 $\mathcal{D}_U$-模。则由推论 \ref{cor:4.5.7}，$i_+(\mathcal{D}(\mathcal{W}))=Ri_\bullet(\mathcal{D}(\mathcal{W}))=\mathcal{D}(i_\bullet(\mathcal{W}))$ 是具有拟凝聚上同调的复形。又因 $\varphi\circ i:U\to Y$ 是仿射态射，由定理 \ref{thm:4.5.3} 有
$\varphi_+(i_+(\mathcal{D}(\mathcal{W})))=(\varphi\circ i)_+(\mathcal{D}(\mathcal{W}))$，而由推论 \ref{cor:4.5.6}，后者是具有拟凝聚上同调的复形。由于 $i_\bullet(\mathcal{W})$ 形式的模生成 $D^b_{qc}(\mathcal{M}_R(\mathcal{D}_X))$，可知该陈述对有界复形成立。最终结论由标准截断点论证得出。
\end{proof}

\section{浸入的直像}

设 $X$ 是光滑代数簇，$Y$ 是其光滑子簇。令 $n=\dim X$、$m=\dim Y$。设 $i:Y\to X$ 是典范浸入。考虑 $X$ 上（相应地 $Y$ 上）右 D-模的范畴 $\mathcal{M}_R(\mathcal{D}_X)$（相应地 $\mathcal{M}_R(\mathcal{D}_Y)$），以及相应的有界导出范畴 $D^b(\mathcal{M}_R(\mathcal{D}_Y))$、$D^b(\mathcal{M}_R(\mathcal{D}_X))$。则有直像函子
$i_+:D^b(\mathcal{M}_R(\mathcal{D}_Y))\to D^b(\mathcal{M}_R(\mathcal{D}_X))$。

\begin{proposition}\label{prop:4.6.1}
模 $\mathcal{D}_{Y\to X}$ 是局部自由 $\mathcal{D}_Y$-模。
\end{proposition}

为证明这个结果，先构造一个适应当前情形的局部坐标系。

\begin{lemma}\label{lem:4.6.2}
设 $Y$ 是光滑 $n$ 维簇 $X$ 的光滑 $m$ 维子簇。设 $y\in Y$。则存在 $y$ 在 $X$ 中的仿射开邻域 $U$ 及 $U$ 上的坐标系
$(f_1,f_2,\dots,f_n;D_1,D_2,\dots,D_n)$，具有如下性质：
\begin{enumerate}[label=(\roman*)]
  \item $U\cap Y$ 是 $U$ 的闭子簇；
  \item $R(U)$ 中所有在 $U\cap Y$ 上消没的函数构成的理想 $I(U\cap Y)$ 由
    $f_{m+1},f_{m+2},\dots,f_n$ 生成；
  \item 向量场 $D_1,D_2,\dots,D_m$ 把理想 $I(U\cap Y)$ 映入自身。
\end{enumerate}
\end{lemma}

\begin{proof}
这个证明是 2.9 证明的一个小变体。由于陈述是局部的，显然可设 $Y$ 在 $X$ 中闭且 $X$ 是某个 $k^p$ 的闭子簇。以 $I_X$、$I_Y$ 分别记 $A=k[X_1,X_2,\dots,X_p]$ 中所有在 $X$、$Y$ 上消没的多项式构成的理想。显然 $I_X\subset I_Y$。由于
$\dim_k T_y(X)=\dim X=n$、$\dim_k T_y(Y)=\dim Y=m$，由 2.6 可找到多项式
$P_{m+1},P_{m+2},\dots,P_p\in I_Y$ 使得：
\begin{enumerate}[label=(\alph*)]
  \item $P_{n+1},P_{n+2},\dots,P_p$ 在 $I_X$ 中；
  \item 矩阵 $[(\partial_iP_j)(y)]$ 的秩为 $p-m$。
\end{enumerate}
这意味着该矩阵的秩在 $y\in k^p$ 的某个邻域 $V$ 上等于 $p-m$，且
$$T_y(X)=\{(\xi_1,\dots,\xi_p)\in k^p\mid
\sum_{i=1}^p\xi_i(\partial_iP_j)(y)=0,\ n+1\le j\le p\},$$
$$T_y(Y)=\{(\xi_1,\dots,\xi_p)\in k^p\mid
\sum_{i=1}^p\xi_i(\partial_iP_j)(y)=0,\ m+1\le j\le p\}.$$
现在如 2.9 的证明，可找到 $g\in A$ 使得 $(I_X)_g$ 在 $A_g$ 中由 $P_{n+1},\dots,P_p$ 生成，$(I_Y)_g$ 由 $P_{m+1},\dots,P_p$ 生成。

可找到多项式 $P_1,P_2,\dots,P_m\in A$ 使得矩阵 $[(\partial_iP_j)(y);1\le i,j\le p]$ 非奇异。因此必要时改变 $g$，还可设它在 $k^p$ 的主开集 $V'$ 上非奇异。以 $Q$ 记此矩阵的逆。则 $Q$ 的矩阵系数在 $A_g$ 中。因此在 $V'$ 上可定义微分算子
$\delta_i=\sum_{j=1}^pQ_{ij}\partial_j$（$1\le i\le n$）。显然对任意 $1\le j\le p$ 满足
$$\delta_i(P_j)=\sum_{k=1}^pQ_{ik}\partial_kP_j=\delta_{ij}.$$
由于任意 $f\in(I_X)_g$ 可表示为 $f=\sum_{j=n+1}^ph_jP_j$（$h_j\in A_g$），有
$$\delta_i(f)=\delta_i\Bigl(\sum_{j=n+1}^ph_jP_j\Bigr)
=\sum_{j=n+1}^p\bigl(\delta_i(h_j)P_j+h_j\delta_i(P_j)\bigr)
=\sum_{j=n+1}^p\delta_i(h_j)P_j\in(I_X)_g,$$
即 $(I_X)_g$ 在 $\delta_i$（$1\le i\le n$）作用下不变。设 $U=X\cap V'$。由于 $R(U)=A_g/(I_X)_g$，这说明 $\delta_i$（$1\le i\le n$）诱导出 $U=X\cap V'$ 上的局部向量场，记为 $D_i$（$1\le i\le n$）。

显然 $I(U\cap Y)=(I_Y)_g/(I_X)_g$，故它由函数 $f_i=P_i|_U$（$m+1\le i\le n$）生成。这证明了 (ii)。

由类似计算还可见 $(I_Y)_g$ 在 $\delta_i$（$1\le i\le m$）作用下不变。因此 $D_1,D_2,\dots,D_m$ 把理想 $I(U\cap Y)$ 映入自身。显然 $D_i(f_j)=\delta_i(P_j)=\delta_{ij}$，故 $(f_1,\dots,f_n;D_1,\dots,D_n)$ 是 $U$ 上的坐标系。
\end{proof}

现在证明命题 \ref{prop:4.6.1}。先设 $X$ 在如下意义下是“小”的。存在 $X$ 上的坐标系 $(f_1,\dots,f_n;D_1,\dots,D_n)$ 使得：
\begin{enumerate}[label=(\roman*)]
  \item 对所有 $I,J\in\mathbb{Z}_+^n$，$D^I\circ D^J=D^{I+J}$；
  \item $(D^I;I\in\mathbb{Z}_+^n)$ 是自由 $R(X)$-模 $\mathcal{D}_X$ 关于左乘（相应地右乘）的一组基；
  \item $R(X)$ 中所有在 $Y$ 上消没的函数构成的理想 $I(Y)$ 由 $f_{m+1},\dots,f_n$ 生成；
  \item 向量场 $D_1,D_2,\dots,D_m$ 把理想 $I(Y)$ 映入自身。
\end{enumerate}
由 3.5 与 8，任意 $y\in Y$ 有这样的“小”邻域 $U$。此外，由 (iv)，$D_1,D_2,\dots,D_m$ 诱导出 $Y$ 上向量场 $T_1,T_2,\dots,T_m$。若以 $g_i$ 记 $f_i$ 在 $Y$ 上的限制（$1\le i\le m$），可见 $(g_1,\dots,g_m;T_1,\dots,T_m)$ 是 $Y$ 上的坐标系。必要时缩小 $X$，由 3.5 还可设
\begin{enumerate}[label=(\roman*),resume]
  \item $(T^I;I\in\mathbb{Z}_+^m)$ 是自由 $R(Y)$-模 $\mathcal{D}_Y$ 关于左乘（相应地右乘）的一组基。
\end{enumerate}
在这些条件下有如下结果。

\begin{lemma}\label{lem:4.6.3}
$(D^I;I\in\{0\}\times\mathbb{Z}_+^{n-m})$ 是自由左 $\mathcal{D}_Y$-模 $\mathcal{D}_{Y\to X}$ 的一组基。
\end{lemma}

\begin{proof}
由 (ii) 可见正合列
$$0\longrightarrow I(Y)\longrightarrow R(X)\longrightarrow R(Y)\longrightarrow 0$$
导致（$R(X)$，右 $\mathcal{D}_X$）-双模的正合列
$$0\longrightarrow I(Y)\otimes_{R(X)}\mathcal{D}_X\longrightarrow\mathcal{D}_X
\longrightarrow\mathcal{D}_{Y\to X}\longrightarrow 0.$$
由于 $D_i$（$1\le i\le m$）保持 $I(Y)$ 不变，左乘 $D_i$ 把 $I(Y)\otimes_{R(X)}\mathcal{D}_X$ 映入自身，并诱导出 $R(Y)$-模 $\mathcal{D}_{Y\to X}$ 的微分自同态，且与 $\mathcal{D}_X$ 的右作用可交换。此外它把 $g\otimes1\in\mathcal{D}_{Y\to X}$ 映为 $T_i(g)\otimes1+g\otimes D_i$，故它等于 $T_i$ 在 $\mathcal{D}_{Y\to X}$ 上的左作用。因此 $T_i$（$1\le i\le m$）在 $\mathcal{D}_{Y\to X}$ 上的左作用为
$$T_i(g\otimes S)=T_i(g)\otimes S+g\otimes D_iS$$
对 $g\in R(Y)$、$S\in\mathcal{D}_X$ 成立。这意味着对任意
$I\in\mathbb{Z}_+^m\times\{0\}\subset\mathbb{Z}_+^n$ 与 $J\in\{0\}\times\mathbb{Z}_+^{n-m}\subset\mathbb{Z}_+^n$ 有
$T^I(1\otimes D^J)=1\otimes D^{I+J}$。特别地，$(D^I;I\in\{0\}\times\mathbb{Z}_+^{n-m})$ 生成左 $\mathcal{D}_Y$-模 $\mathcal{D}_{Y\to X}$。由于由 (ii)，$(D^I;I\in\mathbb{Z}_+^n)$ 是自由 $R(Y)$-模 $\mathcal{D}_{Y\to X}$ 的一组基，由 (v) 可知 $(D^I;I\in\{0\}\times\mathbb{Z}_+^{n-m})$ 是自由左 $\mathcal{D}_Y$-模 $\mathcal{D}_{Y\to X}$ 的一组基。
\end{proof}

因此，由命题 \ref{prop:4.6.1}，对 $D^b(\mathcal{M}_R(\mathcal{D}_Y))$ 中任意有界复形有
$$i_+(\mathcal{V}^\cdot)=Ri_\bullet(\mathcal{V}^\cdot\overset{L}{\otimes}_{\mathcal{D}_Y}\mathcal{D}_{Y\to X})
=Ri_\bullet(\mathcal{V}^\cdot\otimes_{\mathcal{D}_Y}\mathcal{D}_{Y\to X}).$$
下一个结果表明 $i_+$ 对浸入是右导出函子。

\begin{theorem}\label{thm:4.6.4}
设 $Y$ 是光滑簇 $X$ 的光滑子簇，$i:Y\to X$ 是典范浸入。则：
\begin{enumerate}[label=(\roman*)]
  \item 函子 $H^0\circ i_+\circ\mathcal{D}:\mathcal{M}_R(\mathcal{D}_Y)\to\mathcal{M}_R(\mathcal{D}_X)$，由
    $$\mathcal{V}\longmapsto H^0(i_+(\mathcal{D}(\mathcal{V})))
    =i_\bullet(\mathcal{V}\otimes_{\mathcal{D}_Y}\mathcal{D}_{Y\to X})$$
    给出，是左正合的；
  \item 函子 $i_+$ 是 $H^0\circ i_+\circ\mathcal{D}:\mathcal{M}_R(\mathcal{D}_Y)\to\mathcal{M}_R(\mathcal{D}_X)$ 的右导出函子；
  \item $H^0(i_+(\mathcal{D}(\mathcal{V})))$ 的支集等于 $\operatorname{supp}(\mathcal{V})$ 在 $X$ 中的闭包。
\end{enumerate}
\end{theorem}

\begin{proof}
\begin{enumerate}[label=(\roman*)]
  \item 这个断言是显然的，因为
    $$H^0(i_+(\mathcal{D}(\mathcal{V})))=H^0(Ri_\bullet(\mathcal{D}(\mathcal{V})\otimes_{\mathcal{D}_Y}\mathcal{D}_{Y\to X}))$$
    $$=H^0(Ri_\bullet(\mathcal{D}(\mathcal{V}\otimes_{\mathcal{D}_Y}\mathcal{D}_{Y\to X})))
    =i_\bullet(\mathcal{V}\otimes_{\mathcal{D}_Y}\mathcal{D}_{Y\to X}),$$
    而由命题 \ref{prop:4.6.1}，函子 $\mathcal{V}\mapsto\mathcal{V}\otimes_{\mathcal{D}_Y}\mathcal{D}_{Y\to X}$ 是正合的，$i_\bullet$ 是左正合的。
  \item 设 $I$ 是内射右 $\mathcal{D}_Y$-模。存在 $Y$ 的开覆盖 $\{U_j;1\le j\le p\}$ 使得
    $\mathcal{D}_{Y\to X}|_{U_j}$ 是自由 $\mathcal{D}_{U_j}$-模。这意味着
    $(I\otimes_{\mathcal{D}_Y}\mathcal{D}_{Y\to X})|_{U_j}$ 是 $I|_{U_j}$ 无穷多个拷贝的直和。由于
    $I|_{U_j}$ 也是内射的，它是松的（flabby），因此 $(I\otimes_{\mathcal{D}_Y}\mathcal{D}_{Y\to X})|_{U_j}$ 是松的。由此可知正合函子
    $\mathcal{V}\mapsto\mathcal{V}\otimes_{\mathcal{D}_Y}\mathcal{D}_{Y\to X}$ 把 $\mathcal{M}_R(\mathcal{D}_Y)$ 中的内射对象映为松层，即映为对 $i_\bullet$ 无循环的层。这意味着相应导出函子的复合
    $\mathcal{V}^\cdot\mapsto Ri_\bullet(\mathcal{V}^\cdot\otimes_{\mathcal{D}_Y}\mathcal{D}_{Y\to X})$ 是左正合函子
    $\mathcal{V}\mapsto i_\bullet(\mathcal{V}\otimes_{\mathcal{D}_Y}\mathcal{D}_{Y\to X})$ 的右导出函子。
  \item 设 $U$ 是 $X$ 中开集。则
    $$H^0(i_+(\mathcal{D}(\mathcal{V})))(U)=i_\bullet(\mathcal{V}\otimes_{\mathcal{D}_Y}\mathcal{D}_{Y\to X})(U)
    =(\mathcal{V}\otimes_{\mathcal{D}_Y}\mathcal{D}_{Y\to X})(U\cap Y).$$
    因此 $H^0(i_+(\mathcal{D}(\mathcal{V})))$ 在 $\operatorname{supp}(\mathcal{V})$ 的补集的任一开集上的限制等于 0。故
    $$\operatorname{supp}(H^0(i_+(\mathcal{D}(\mathcal{V}))))\subset\operatorname{supp}(\mathcal{V}).$$
    设 $y\in\operatorname{supp}(\mathcal{V})$。先假设 $y$ 不在 $\operatorname{supp}(H^0(i_+(\mathcal{D}(\mathcal{V}))))$ 中。则存在引理 \ref{lem:4.6.2} 所述的 $y$ 在 $X$ 中的“小”仿射邻域 $U$，使得 $H^0(i_+(\mathcal{D}(\mathcal{V})))(U)=0$。由引理 \ref{lem:4.6.3}，$(D^I;I\in\{0\}\times\mathbb{Z}_+^{n-m})$ 是自由 $\mathcal{D}_Y$-模 $\mathcal{D}_{Y\to X}$ 的一组基。这说明 $(\mathcal{V}\otimes_{\mathcal{D}_Y}\mathcal{D}_{Y\to X})|_{U\cap Y}$ 是 $\mathcal{V}|_{U\cap Y}$ 无穷多个拷贝的直和。故 $H^0(i_+(\mathcal{D}(\mathcal{V})))(U)=0$ 将蕴含 $\mathcal{V}|_{U\cap Y}=0$，这是不可能的。因此 $y\in\operatorname{supp}(H^0(i_+(\mathcal{D}(\mathcal{V}))))$。于是
    $$\operatorname{supp}(\mathcal{V})\subset\operatorname{supp}(H^0(i_+(\mathcal{D}(\mathcal{V})))).$$
    这证明了等式。
\end{enumerate}
\end{proof}

结合前述结果，得到以下推论。

\begin{corollary}\label{cor:4.6.5}
设 $i:Y\to X$ 是仿射浸入。则函子
$H^0\circ i_+\circ\mathcal{D}:\mathcal{M}^R_{qc}(\mathcal{D}_Y)\to\mathcal{M}^R_{qc}(\mathcal{D}_X)$ 是正合的。
\end{corollary}

若 $Y$ 是 $X$ 的闭光滑子簇，$i_\bullet$ 是正合函子。因此有如下结果。

\begin{corollary}\label{cor:4.6.6}
设 $Y$ 是 $X$ 的闭光滑子簇。则函子
$H^0\circ i_+\circ\mathcal{D}:\mathcal{M}_R(\mathcal{D}_Y)\to\mathcal{M}_R(\mathcal{D}_X)$ 是正合的。
\end{corollary}

\begin{proof}
由于 $i_\bullet$ 是正合函子，
$$\mathcal{V}\longmapsto H^0(i_+(\mathcal{D}(\mathcal{V})))=i_\bullet(\mathcal{V}\otimes_{\mathcal{D}_Y}\mathcal{D}_{Y\to X})$$
是从 $\mathcal{M}_R(\mathcal{D}_Y)$ 到 $\mathcal{M}_R(\mathcal{D}_X)$ 的正合函子。
\end{proof}

借由记号上的滥用，此时也把函子 $H^0\circ i_+\circ\mathcal{D}$ 记为 $i_+$。因此有
$$i_+(\mathcal{V})=i_\bullet(\mathcal{V}\otimes_{\mathcal{D}_Y}\mathcal{D}_{Y\to X})$$
对任意 $\mathcal{M}_R(\mathcal{D}_Y)$ 中的 $\mathcal{V}$ 成立，且对任意 $D^b(\mathcal{M}_R(\mathcal{D}_X))$ 中的 $\mathcal{V}^\cdot$ 有
$(i_+(\mathcal{V}^\cdot))^p=i_+(\mathcal{V}^p)$。

\begin{proposition}\label{prop:4.6.7}
设 $Y$ 是 $X$ 的闭光滑子簇。则对任意凝聚右 $\mathcal{D}_Y$-模 $\mathcal{V}$，右 $\mathcal{D}_X$-模 $i_+(\mathcal{V})$ 也是凝聚的。
\end{proposition}

\begin{proof}
只需证明对 $X$ 中任意仿射开集 $U$，$i_+(\mathcal{V})|_U$ 是凝聚 $\mathcal{D}_U$-模。因此用 $U$ 代替 $X$、$Y\cap U$ 代替 $Y$，可设 $X$ 是光滑仿射簇。此时 $i_+(\mathcal{V})$ 是右 $\mathcal{D}_X$-模 $\Gamma(Y,\mathcal{V})\otimes_{\mathcal{D}_Y}\mathcal{D}_{Y\to X}$ 的局部化。由前述，只需证 $\Gamma(Y,\mathcal{V})\otimes_{\mathcal{D}_Y}\mathcal{D}_{Y\to X}$ 是有限生成右 $\mathcal{D}_X$-模。由于 $Y$ 是 $X$ 的闭子簇，$R(Y)$ 是 $R(X)$ 的商，且 $\mathcal{D}_{Y\to X}=R(Y)\otimes_{R(X)}\mathcal{D}_X$ 是 $\mathcal{D}_X$ 作为右 $\mathcal{D}_X$-模的商。因此它由元素 $1\otimes1$ 生成。由假设，$\Gamma(Y,\mathcal{V})$ 是有限生成 $\mathcal{D}_Y$-模。设 $v_1,\dots,v_n$ 是它的一组生成元。则 $v_1\otimes1\otimes1,\dots,v_n\otimes1\otimes1$ 生成
$\Gamma(Y,\mathcal{V})\otimes_{\mathcal{D}_Y}\mathcal{D}_{Y\to X}$ 作为右 $\mathcal{D}_X$-模。
\end{proof}

特别地，若 $\mathcal{V}$ 是凝聚右 $\mathcal{D}_Y$-模，可比较 $\mathcal{V}$ 与 $i_+(\mathcal{V})$ 的特征簇。设 $x\in Y$。则 $T_x(Y)\subset T_x(X)$，故有自然投影 $p_x:T_x^*(X)\to T_x^*(Y)$。

设 $Y$ 是代数簇 $X$ 的闭子集。定义 $\dim_YX=\sup_{x\in Z}(\dim_xX)$（原文如此，应为 $x\in Y$）。

\begin{theorem}\label{thm:4.6.8}
设 $Y$ 是 $X$ 的闭光滑子簇。则对任意凝聚右 $\mathcal{D}_Y$-模 $\mathcal{V}$ 有
$$\operatorname{Ch}(i_+(\mathcal{V}))=\{(x,\omega)\in T^*(X)\mid(x,p_x(\omega))\in\operatorname{Ch}(\mathcal{V})\}.$$
特别地，
$$\dim_{p_X^{-1}(x)}\operatorname{Ch}(i_+(\mathcal{V}))
=\dim_{p_Y^{-1}(x)}\operatorname{Ch}(\mathcal{V})+\operatorname{codim}_xY.$$
\end{theorem}

\begin{proof}
由于 $i_+(\mathcal{V})$ 的支集在 $Y$ 中，只需证明每个 $y\in Y$ 有 $X$ 中开邻域 $U$ 使得
$$\operatorname{Ch}(i_+(\mathcal{V})|_U)=\{(x,\omega)\in T^*(U)\mid(x,p_x(\omega))\in\operatorname{Ch}(\mathcal{V}|_{Y\cap U})\}.$$
因此可设该邻域满足引理 \ref{lem:4.6.2} 的条件。用此邻域代替 $X$，可设 $X$ 是“小”的。此时由引理 \ref{lem:4.6.3}，$\mathcal{D}_{Y\to X}$ 是自由左 $\mathcal{D}_Y$-模。设 $V=\Gamma(Y,\mathcal{V})$。于是 $i_+(\mathcal{V})$ 是右 $\mathcal{D}_X$-模 $V\otimes_{\mathcal{D}_Y}\mathcal{D}_{Y\to X}$ 的局部化，且作为向量空间
$$\Gamma(X,i_+(\mathcal{V}))=V\otimes_{\mathcal{D}_Y}\mathcal{D}_{Y\to X}
=\bigoplus_{I\in\{0\}\times\mathbb{Z}_+^{n-m}}V\otimes_kD^I.$$
设 $FV$ 是右 $\mathcal{D}_Y$-模 $V$ 的好滤过。则由 I.3.1，$\operatorname{Gr}V$ 是有限生成 $\operatorname{Gr}\mathcal{D}_Y$-模。可通过下式定义 $\Gamma(X,i_+(\mathcal{V}))$ 的滤过：
$$F_p\Gamma(X,i_+(\mathcal{V}))=\bigoplus_{I\in\{0\}\times\mathbb{Z}_+^{n-m},\ s+|I|\le p}F_sV\otimes D^I.$$
显然这是穷尽的递增滤过，且对充分负的 $p\in\mathbb{Z}$ 有 $F_p\Gamma(X,i_+(\mathcal{V}))=0$。

断言这是右 $\mathcal{D}_X$-模滤过。设 $v\otimes D^I\in\Gamma(X,i_+(\mathcal{V}))$，
$v\in F_sV$，$I\in\{0\}\times\mathbb{Z}_+^{n-m}$。断言对任意 $f\in R(X)$ 有
$$(v\otimes D^I)f\in F_{s+|I|}\Gamma(X,i_+(\mathcal{V})).$$
对 $|I|$ 作归纳。若 $|I|=0$，则
$$(v\otimes1)f=vf|_Y\otimes1\in F_sV\otimes1.$$
设 $|I|>0$。则可找到 $m+1\le j\le n$ 与 $I'\in\{0\}\times\mathbb{Z}_+^{n-m}$、$|I'|=|I|-1$，使得 $D^I=D^{I'}D_j$。故
$$(v\otimes D^I)f=(v\otimes D^{I'})D_jf=(v\otimes D^{I'})D_j(f)+(v\otimes D^{I'})fD_j.$$
由归纳假设，$(v\otimes D^{I'})D_j(f)$ 与 $(v\otimes D^{I'})f$ 都在 $F_{s+|I'|}\Gamma(X,i_+(\mathcal{V}))$ 中。故
$$(v\otimes D^{I'})fD_j\in F_{s+|I'|}\Gamma(X,i_+(\mathcal{V}))D_j
\subset F_{s+|I|}\Gamma(X,i_+(\mathcal{V})).$$
这证明了我们的断言，即 $F\Gamma(X,i_+(\mathcal{V}))$ 是由 $R(X)$-子模构成的滤过。

设 $1\le j\le n$。断言有
$$(v\otimes D^I)D_j\in F_{s+|I|+1}\Gamma(X,i_+(\mathcal{V})).$$
当 $m+1\le j\le n$ 时显然。若 $1\le j\le m$，则
$$(v\otimes D^I)D_j=vT_j\otimes D^I\in F_{s+|I|+1}\Gamma(X,i_+(\mathcal{V})),$$
因为 $vT_j\in F_{s+1}V$。由前述，这说明 $F\Gamma(X,i_+(\mathcal{V}))$ 是 $\mathcal{D}_X$-模滤过。

显然
$$\operatorname{Gr}\Gamma(X,i_+(\mathcal{V}))=\operatorname{Gr}V\otimes_kk[\xi_{m+1},\dots,\xi_n],$$
且 $\operatorname{Gr}_p\Gamma(X,i_+(\mathcal{V}))$ 由 $v\otimes\xi^I$（$v\in\operatorname{Gr}_sV$，$I\in\{0\}\times\mathbb{Z}_+^{n-m}$，$|I|=p-s$）张成。又有
$\operatorname{Gr}\mathcal{D}_X=R(X)[\xi_1,\dots,\xi_n]=R(X)[\xi_1,\dots,\xi_m]\otimes_kk[\xi_{m+1},\dots,\xi_n]$。$\operatorname{Gr}\mathcal{D}_X$ 在 $\operatorname{Gr}\Gamma(X,i_+(\mathcal{V}))$ 上的作用如下：$f\in R(X)$ 乘以第一因子中限制 $f|_Y$；$\xi_j$（$1\le j\le m$）在第一因子中相乘；$\xi_j$（$m+1\le j\le n$）在第二因子中相乘。

因此，若 $v_1,\dots,v_k$ 是 $\operatorname{Gr}V$ 作为 $\operatorname{Gr}\mathcal{D}_Y$-模的生成元，则 $v_1\otimes1,\dots,v_k\otimes1$ 生成 $\operatorname{Gr}\Gamma(X,i_+(\mathcal{V}))$ 作为 $\operatorname{Gr}\mathcal{D}_X$-模。于是 $\operatorname{Gr}\Gamma(X,i_+(\mathcal{V}))$ 是有限生成 $\operatorname{Gr}\mathcal{D}_X$-模。故由 I.3.1，$F\Gamma(X,i_+(\mathcal{V}))$ 是好滤过。

设 $A\subset R(Y)[\xi_1,\dots,\xi_m]$ 是 $\operatorname{Gr}V$ 在 $\operatorname{Gr}\mathcal{D}_Y$ 中的零化子。限制映射 $f\mapsto f|_Y$ 定义了 $R(X)$ 到 $R(Y)$ 的满同态。它诱导 $R(X)[\xi_1,\dots,\xi_m]$ 到 $R(Y)[\xi_1,\dots,\xi_m]$ 的满同态 $\varphi$。令 $B=\varphi^{-1}(A)$。则 $B\otimes_kk[\xi_{m+1},\dots,\xi_n]$ 零化 $\operatorname{Gr}\Gamma(X,i_+(\mathcal{V}))$。设 $S\otimes\xi^J$ 是 $\operatorname{Gr}\Gamma(X,i_+(\mathcal{V}))$ 零化子中的齐次元素。则对 $v\in\operatorname{Gr}V$ 有
$(S\otimes\xi^J)(v\otimes1)=Sv\otimes\xi^J=0$，即 $Sv=0$，$S\in B$。这说明 $\operatorname{Gr}\Gamma(X,i_+(\mathcal{V}))$ 的零化子等于 $B\otimes_kk[\xi_{m+1},\dots,\xi_n]$。

我们把 $R(T^*(Y))$ 与 $R(Y)[\xi_1,\dots,\xi_m]$ 的等同，对应于把 $T^*(Y)$ 与 $Y\times k^m$ 通过
$(y,\eta)\mapsto(y,(\eta(T_1(y)),\dots,\eta(T_m(y))))$ 的等同。类似地，$R(T^*(X))$ 与 $R(X)[\xi_1,\dots,\xi_n]$ 的等同对应于把 $T^*(X)$ 与 $X\times k^n$ 通过 $(x,\omega)\mapsto(x,(\omega(D_1(x)),\dots,\omega(D_n(x))))$ 的等同。在这些等同下，特征簇 $\operatorname{Ch}(\mathcal{V})$ 对应于 $Y\times k^m$ 中 $A$ 的零点集构成的子簇，而 $\operatorname{Ch}(i_+(\mathcal{V}))$ 对应于 $X\times k^n$ 中
$B\otimes_kk[\xi_{m+1},\dots,\xi_n]$ 的零点集构成的子簇。可通过映射
$(y,\xi_1,\dots,\xi_m)\mapsto(y,\xi_1,\dots,\xi_m,0,\dots,0)$ 把 $Y\times k^m$ 嵌入 $X\times k^n$。这个嵌入对应于 $R(X)[\xi_1,\dots,\xi_n]$ 到 $R(Y)[\xi_1,\dots,\xi_m]$ 的自然投影。在此嵌入下有等同 $\operatorname{Ch}(i_+(\mathcal{V}))=\operatorname{Ch}(\mathcal{V})\times k^{n-m}$。这立即蕴含
$$\dim\operatorname{Ch}(i_+(\mathcal{V}))=\dim\operatorname{Ch}(\mathcal{V})+n-m.$$
此外，点 $(x,\omega)$ 属于 $\operatorname{Ch}(i_+(\mathcal{V}))$ 当且仅当
$(x,(\omega(D_1(x)),\dots,\omega(D_m(x))))$ 对应于 $\operatorname{Ch}(\mathcal{V})$ 中的点，故 $x\in Y$。由于对 $x\in Y$、$1\le i\le m$ 有 $\omega(D_i(x))=p_x(\omega)(T_i(x))$，可见 $(x,\omega)\in\operatorname{Ch}(i_+(\mathcal{V}))$ 等价于 $x\in Y$ 且 $(x,p_x(\omega))\in\operatorname{Ch}(\mathcal{V})$。
\end{proof}

设 $\mathcal{V}$ 是 $X$ 上非零凝聚 $\mathcal{D}_X$-模。对任意 $x\in X$，记
$$\operatorname{holdef}_x(\mathcal{V})=\dim_{p_X^{-1}(x)}\operatorname{Ch}(\mathcal{V})-\dim_xX.$$
称这个数为 $\mathcal{V}$ 在 $x$ 处的\emph{完整亏格（holonomic defect）}。显然对任意 $x\notin\operatorname{supp}(\mathcal{V})$，$\mathcal{V}$ 的完整亏格等于 $-\dim_xX$。

\begin{corollary}\label{cor:4.6.9}
设 $Y$ 是 $X$ 的闭光滑子簇。则对任意非零凝聚右 $\mathcal{D}_Y$-模 $\mathcal{V}$ 与任意 $x\in Y$ 有
$$\operatorname{holdef}_x(\mathcal{V})=\operatorname{holdef}_x(i_+(\mathcal{V})).$$
\end{corollary}

\section{Bernstein 不等式}

设 $X$ 是光滑簇。下一个结果对 D-模理论具有根本重要性。它推广了前述结论。

\begin{theorem}[Bernstein 不等式]\label{thm:4.7.1}
设 $\mathcal{V}$ 是凝聚 $\mathcal{D}_X$-模。则对任意 $x\in\operatorname{supp}(\mathcal{V})$ 有
$$\dim_{p_X^{-1}(x)}\operatorname{Ch}(\mathcal{V})\ge\dim_xX.$$
\end{theorem}

\begin{proof}
设 $\mathcal{V}\ne0$。则可取一点 $x\in\operatorname{supp}(\mathcal{V})$ 及 $x$ 的连通开仿射邻域 $U$，使得
$p_X^{-1}(U)$ 只与 $\operatorname{Ch}(\mathcal{V})$ 中和 $p_X^{-1}(x)$ 相交的不可约分支相交。因此
$\dim U=\dim_xX$，且 $\mathcal{V}|_U$ 是非零凝聚 $\mathcal{D}_U$-模，满足
$\dim\operatorname{Ch}(\mathcal{V}|_U)=\dim_{p_X^{-1}(x)}\operatorname{Ch}(\mathcal{V})$。
显然只需对 $\mathcal{V}|_U$ 证明该陈述。此时可设 $U$ 是某个 $k^n$（$n\in\mathbb{Z}_+$）的闭光滑子簇。设
$i:U\to k^n$ 是相应的闭浸入。由前述，
$$\dim_{p_X^{-1}(x)}\operatorname{Ch}(\mathcal{V})-\dim_xX
=\dim_{p_U^{-1}(x)}\operatorname{Ch}(\mathcal{V}|_U)-\dim_xU=\operatorname{holdef}_x(\mathcal{V}|_U)$$
$$=\operatorname{holdef}_x(i_+(\mathcal{V}|_U))=\dim\operatorname{Ch}(i_+(\mathcal{V}|_U))-n.$$
而由前述可知 $\dim\operatorname{Ch}(i_+(\mathcal{V}|_U))\ge n$。
\end{proof}
\begin{annote}
Bernstein 不等式 $\dim_{\pi^{-1}(x)}\mathrm{Ch}(\mathcal V)\ge\dim_x X$ 的证明是“化到局部”：取 $x$ 的仿射邻域并令 $f$ 为到该邻域的光滑投影，用 $i_+$ 把问题搬到特征簇上，再用完整亏格 $\mathrm{holdef}$ 在 $i_+$ 下的不变性把两项维数对齐。讲义此处的 $\mathrm{holdef}$ 链条与原文一致；注意第二步的 $\dim\mathrm{Ch}(i_+(\mathcal V|_U))\ge n$ 在原书里写作“$\mathrm{Ch}(i_+(\mathcal V|_U))\ge n$”，量纲不严，译文照原文。[来源：原书 §7, p119–124；nLab 条目 holonomic+D-module（2014-02-16）]
\end{annote}



\section{闭浸入与 Kashiwara 定理}

设 $X$ 是光滑簇，$Y$ 是其闭光滑子簇。设 $i:Y\to X$ 是 $Y$ 到 $X$ 的自然包含。我们在前述中证明了直像函子
$$i_+(\mathcal{V})=i_\bullet(\mathcal{V}\otimes_{\mathcal{D}_Y}\mathcal{D}_{Y\to X})$$
是从 $\mathcal{M}_R(\mathcal{D}_Y)$ 到 $\mathcal{M}_R(\mathcal{D}_X)$ 的正合函子。

\begin{proposition}\label{prop:4.8.1}
函子 $i_+:\mathcal{M}_R(\mathcal{D}_Y)\to\mathcal{M}_R(\mathcal{D}_X)$ 有右伴随
$i^!:\mathcal{M}_R(\mathcal{D}_X)\to\mathcal{M}_R(\mathcal{D}_Y)$。
\end{proposition}

\begin{proof}
显然，由于 $i:Y\to X$ 是闭浸入，直像函子
$i_\bullet:\mathcal{M}_R(i^{-1}(\mathcal{D}_X))\to\mathcal{M}_R(\mathcal{D}_X)$ 有右伴随
$i^{-1}:\mathcal{M}_R(\mathcal{D}_X)\to\mathcal{M}_R(i^{-1}(\mathcal{D}_X))$。因此对
$\mathcal{M}_R(\mathcal{D}_Y)$ 与 $\mathcal{M}_R(\mathcal{D}_X)$ 中的 $\mathcal{V}$、$\mathcal{W}$ 有
$$\operatorname{Hom}_{\mathcal{D}_X}(i_+(\mathcal{V}),\mathcal{W})
=\operatorname{Hom}_{\mathcal{D}_X}(i_\bullet(\mathcal{V}\otimes_{\mathcal{D}_Y}\mathcal{D}_{Y\to X}),\mathcal{W})$$
$$=\operatorname{Hom}_{i^{-1}(\mathcal{D}_X)}(\mathcal{V}\otimes_{\mathcal{D}_Y}\mathcal{D}_{Y\to X},i^{-1}(\mathcal{W})).$$
现在利用张量积的性质，得
$$\operatorname{Hom}_{\mathcal{D}_X}(i_+(\mathcal{V}),\mathcal{W})
=\operatorname{Hom}_{\mathcal{D}_Y}(\mathcal{V},
\operatorname{Hom}_{i^{-1}(\mathcal{D}_X)}(\mathcal{D}_{Y\to X},i^{-1}(\mathcal{W}))).$$
因此函子
$$i^!(\mathcal{W})=\operatorname{Hom}_{i^{-1}(\mathcal{D}_X)}(\mathcal{D}_{Y\to X},i^{-1}(\mathcal{W}))$$
是 $i_+$ 的右伴随。
\end{proof}

显然，右伴随 $i^!:\mathcal{M}_R(\mathcal{D}_X)\to\mathcal{M}_R(\mathcal{D}_Y)$ 是左正合函子。此外，由于 $i_+$ 正合，$i^!$ 把内射模映为内射模。

现在想找 $i^!$ 的另一种描述。显然
$$i^!(\mathcal{W})=\operatorname{Hom}_{i^{-1}(\mathcal{D}_X)}(\mathcal{D}_{Y\to X},i^{-1}(\mathcal{W}))$$
$$=\operatorname{Hom}_{i^{-1}(\mathcal{D}_X)}(\mathcal{O}_Y\otimes_{i^{-1}(\mathcal{O}_X)}i^{-1}(\mathcal{D}_X),i^{-1}(\mathcal{W}))
=\operatorname{Hom}_{i^{-1}(\mathcal{O}_X)}(\mathcal{O}_Y,i^{-1}(\mathcal{W})).$$
设 $\mathcal{J}$ 是 $\mathcal{O}_X$ 中由沿 $Y$ 消没的函数构成的理想层。则有自然正合列
$$0\longrightarrow i^{-1}(\mathcal{J})\longrightarrow i^{-1}(\mathcal{O}_X)\longrightarrow\mathcal{O}_Y\longrightarrow0,$$
它导致正合列
$$0\longrightarrow\operatorname{Hom}_{i^{-1}(\mathcal{O}_X)}(\mathcal{O}_Y,i^{-1}(\mathcal{W}))
\longrightarrow\operatorname{Hom}_{i^{-1}(\mathcal{O}_X)}(i^{-1}(\mathcal{O}_X),i^{-1}(\mathcal{W}))$$
$$\longrightarrow\operatorname{Hom}_{i^{-1}(\mathcal{O}_X)}(i^{-1}(\mathcal{J}),i^{-1}(\mathcal{W})).$$
映射 $T\mapsto T(1)$ 把 $\operatorname{Hom}_{i^{-1}(\mathcal{O}_X)}(i^{-1}(\mathcal{O}_X),i^{-1}(\mathcal{W}))$ 与
$i^{-1}(\mathcal{W})$ 等同。在此同构下，$\operatorname{Hom}_{i^{-1}(\mathcal{O}_X)}(\mathcal{O}_Y,i^{-1}(\mathcal{W}))$ 对应于 $i^{-1}(\mathcal{W})$ 中所有被 $i^{-1}(\mathcal{J})$ 的所有元素零化的截面构成的子层 $\mathcal{W}_0$。

设 $Z$ 是 $Y$ 的闭光滑子簇，$j:Z\to Y$ 是自然包含。则由前述 $(i\circ j)_+=i_+\circ j_+$。因此由伴随函子的唯一性得出以下事实。

\begin{lemma}\label{lem:4.8.2}
有 $(i\circ j)^!\cong j^!\circ i^!$。
\end{lemma}

最后，函子 $i^!$ 保持拟凝聚性，即有如下结果。

\begin{proposition}\label{prop:4.8.3}
设 $\mathcal{V}$ 是拟凝聚右 $\mathcal{D}_X$-模。则 $i^!(\mathcal{V})$ 是拟凝聚右 $\mathcal{D}_Y$-模。
\end{proposition}

\begin{proof}
为证明这个结果，首先注意这个结果是局部的，且只需考虑 $\mathcal{O}_Y$-模结构。因此可设 $X$ 是连通仿射簇，带整体坐标系 $(f_1,\dots,f_n;D_1,\dots,D_n)$，使得 $Y$ 是 $(f_{m+1},\dots,f_n)$ 的零点集。此时集合 $Y_k$（$m\le k\le n$）是 $(f_{k+1},\dots,f_n)$ 的公共零点集，是包含 $Y$ 的 $X$ 的闭光滑子簇。由于 $Y_m=Y$ 且 $Y_m\subset Y_{m+1}$ 余维为 1，由引理 \ref{lem:4.8.2} 可把证明归结为 $Y$ 在 $X$ 中余维为 1 的情形。因此可设 $Y$ 是函数 $f=f_n$ 的零点集。于是由于此情形下 $\mathcal{J}$ 由 $f$ 生成，可考虑正合列
$$0\longrightarrow\mathcal{O}_X\xrightarrow{\ f\ }\mathcal{O}_X\longrightarrow\mathcal{O}_X/\mathcal{J}\longrightarrow0,$$
其中第一个态射是乘以 $f$。把它限制到 $Y$，得到正合列
$$0\longrightarrow i^{-1}(\mathcal{O}_X)\xrightarrow{\ f\ }i^{-1}(\mathcal{O}_X)\longrightarrow\mathcal{O}_Y\longrightarrow0.$$
因此复形
$$0\longrightarrow i^{-1}(\mathcal{O}_X)\xrightarrow{\ f\ }i^{-1}(\mathcal{O}_X)\longrightarrow0$$
是 $\mathcal{O}_Y$ 关于 $i^{-1}(\mathcal{O}_X)$-模的自由分解。把这个分解在 $i^{-1}(\mathcal{D}_X)$ 上与 $i^{-1}(\mathcal{V})$ 作张量积，得到复形
$$0\longrightarrow i^{-1}(\mathcal{V})\xrightarrow{\ f\ }i^{-1}(\mathcal{V})\longrightarrow0,$$
它表示 $Li^*(\mathcal{D}(\mathcal{V}))$。由此可知，作为 $\mathcal{O}_Y$-模有 $i^!(\mathcal{V})=L^{-1}i^*(\mathcal{V})$。由前述，这意味着对拟凝聚 $\mathcal{D}_X$-模 $\mathcal{V}$，$\mathcal{O}_Y$-模 $i^!(\mathcal{V})$ 也是拟凝聚的。
\end{proof}

现在研究伴随态射 $\operatorname{id}\to i^!\circ i_+$ 与 $i_+\circ i^!\to\operatorname{id}$。

由于 $\mathcal{D}_{Y\to X}=\mathcal{O}_Y\otimes_{i^{-1}(\mathcal{O}_X)}i^{-1}(\mathcal{D}_X)$，它由 $1\otimes1$ 决定一个典范截面。因此有由
$T\mapsto T(1\otimes1)$ 给出的典范态射 $\mathcal{D}_Y\to\mathcal{D}_{Y\to X}$。由前述 $\mathcal{D}_{Y\to X}$ 作为 $\mathcal{D}_Y$-模的局部描述可知，这个态射是单射，且有正合列
$$0\longrightarrow\mathcal{D}_Y\longrightarrow\mathcal{D}_{Y\to X}\longrightarrow Q\longrightarrow0$$
其中 $Q$ 是局部自由 $\mathcal{D}_Y$-模。此外，$\mathcal{D}_Y$ 的像包含在 $\mathcal{D}_{Y\to X}$ 中被 $i^{-1}(\mathcal{J})$ 的所有元素零化的截面构成的子层中。

设 $\mathcal{V}$ 是右 $\mathcal{D}_Y$-模。则在 $\mathcal{D}_Y$ 上与典范态射 $\mathcal{D}_Y\to\mathcal{D}_{Y\to X}$ 作张量积，得到自然单射
$$\mathcal{V}\longrightarrow\mathcal{V}\otimes_{\mathcal{D}_Y}\mathcal{D}_{Y\to X}=i^{-1}(i_+(\mathcal{V})),$$
其像包含在 $i^!(i_+(\mathcal{V}))=(i_+(\mathcal{V}))_0$ 中。因此有典范单射 $\alpha_{\mathcal{V}}:\mathcal{V}\to i^!(i_+(\mathcal{V}))$。显然这个态射正是伴随态射之一。

设 $y\in Y$。则可找到 $y$ 的连通仿射邻域，带坐标系 $(f_1,\dots,f_n;D_1,\dots,D_n)$，使得 $Y$ 是 $f_{m+1},\dots,f_n$ 的零点集。此时如前述，$\mathcal{D}_{Y\to X}|_U$ 是以 $(D^I;I\in\{0\}\times\mathbb{Z}_+^{n-m})$ 为基的自由 $\mathcal{D}_{Y\cap U}$-模。令
$$E=\sum_{j=m+1}^nf_jD_j.$$
以下结果由直接计算得出。

\begin{lemma}\label{lem:4.8.4}
\begin{enumerate}[label=(\roman*)]
  \item $[E,f_j]=f_j$，$m+1\le j\le n$；
  \item $[E,D^I]=-|I|D^I$，$I\in\{0\}\times\mathbb{Z}_+^{n-m}$。
\end{enumerate}
\end{lemma}

设 $\mathcal{V}$ 是右 $\mathcal{D}_Y$-模。则
$$i_+(\mathcal{V})_y=\mathcal{V}_y\otimes_{\mathcal{D}_{Y,y}}\mathcal{D}_{Y\to X,y}$$
且 $\mathcal{D}_{Y\to X,y}$ 是自由左 $\mathcal{D}_{Y,y}$-模。此外，$D^I$（$I\in\{0\}\times\mathbb{Z}_+^{n-m}$）的像构成自由 $\mathcal{D}_{Y,y}$-模 $\mathcal{D}_{Y\to X,y}$ 的一组基。故有
$$i_+(\mathcal{V})_y=\bigoplus_{I\in\{0\}\times\mathbb{Z}_+^{n-m}}\mathcal{V}_yD^I.$$
设 $v\in i_+(\mathcal{V})_y$。则有唯一分解 $v=\sum_Iv_ID^I$，$v_I\in\mathcal{V}_y$。定义
$$\operatorname{ord}(v)=\max\{|I|\mid v_I\ne0\}.$$
由于 $\sum_Iv_IE=\sum_jv_If_jD_j=0$，有
$$vE=\sum_Iv_ID^IE=\sum_Iv_I[D^I,E]=\sum_I|I|v_ID^I.$$
因此有如下结果。

\begin{lemma}\label{lem:4.8.5}
设 $v\in i_+(\mathcal{V})_y$。则
\begin{enumerate}[label=(\roman*)]
  \item $\operatorname{ord}(vf_j)\le\operatorname{ord}(v)-1$，$m+1\le j\le n$；
  \item $\operatorname{ord}\bigl(v(E-\operatorname{ord}(v))\bigr)\le\operatorname{ord}(v)-1$。
\end{enumerate}
\end{lemma}

\begin{proof}
\begin{enumerate}[label=(\roman*)]
  \item 设 $v=\sum_Iv_ID^I$，$v_I\in\mathcal{V}_y$。则
    $$vf_j=\sum_Iv_ID^If_j=\sum_Iv_I[D^I,f_j].$$
    若 $I=(i_1,\dots,i_j,\dots,i_n)$ 且 $I'=(i_1,\dots,i_j-1,\dots,i_n)$，可见
    $[D^I,f_j]=i_jD^{I'}$。这立即使 (i) 成立。
  \item 由上述公式立即得出。
\end{enumerate}
\end{proof}

现在可以分析第一个伴随态射。

\begin{lemma}\label{lem:4.8.6}
态射 $\alpha_{\mathcal{V}}$ 诱导 $\mathcal{V}$ 到 $i^!(i_+(\mathcal{V}))$ 的同构。
\end{lemma}

\begin{proof}
只需证明 $\alpha_{\mathcal{V}}$ 是满射，或等价地，$\alpha_{\mathcal{V},y}:\mathcal{V}_y\to i^!(i_+(\mathcal{V}))_y$ 是满射。设 $v$ 是 $i^!(i_+(\mathcal{V}))_y$ 中截面的芽。则对 $m+1\le j\le n$ 有 $vf_j=0$。故 $vE=0$，由上述计算得 $\operatorname{ord}(v)=0$，即 $v\in\mathcal{V}_y$。
\end{proof}

因此，伴随态射 $\operatorname{id}\to i^!\circ i_+$ 是函子的同构。

现在研究另一个伴随态射 $i_+\circ i^!\to\operatorname{id}$。设 $\mathcal{W}$ 是 $\mathcal{D}_X$-模，$\beta_{\mathcal{W}}:i_+(i^!(\mathcal{W}))\to\mathcal{W}$ 是相应自然态射。

\begin{lemma}\label{lem:4.8.7}
态射 $\beta_{\mathcal{W}}:i_+(i^!(\mathcal{W}))\to\mathcal{W}$ 是单射。
\end{lemma}

\begin{proof}
设 $y\in Y$。须证 $\beta_{\mathcal{W},y}:i_+(i^!(\mathcal{W}))_y\to\mathcal{W}_y$ 是单射。设 $\beta_{\mathcal{W},y}$ 不是单射。取 $\beta_{\mathcal{W},y}$ 的核中元素 $v\ne0$。若 $p=\operatorname{ord}(v)>0$，则存在 $m+1\le j\le n$ 使 $vf_j\ne0$。由引理 \ref{lem:4.8.5}(i)，$\operatorname{ord}(vf_j)\le p-1$。故由对 $\operatorname{ord}(v)$ 的递减归纳，得出 $\beta_{\mathcal{W},y}$ 的核中存在 $\operatorname{ord}(v)=0$ 的元素 $v\ne0$。但这显然不可能，因为 $i^!(\mathcal{W})_y$ 是 $\mathcal{W}_y$ 的子空间。
\end{proof}

因此可把 $i_+(i^!(\mathcal{W}))$ 看作 $\mathcal{W}$ 的子模。由于它支集在 $Y$ 中，它是 $\mathcal{W}$ 中所有支集在 $Y$ 中的局部截面构成的模 $\Gamma_Y(\mathcal{W})$ 的子模。

归纳构造一列递减的理想层 $(\mathcal{J}^p;p\in\mathbb{Z}_+)$：
$$\mathcal{J}^0=\mathcal{O}_X,\qquad \mathcal{J}^p\text{ 是 }\mathcal{J}^{p-1}\otimes_{\mathcal{O}_X}\mathcal{J}\to\mathcal{O}_X\text{ 的像},\quad p\in\mathbb{N}.$$
设 $\mathcal{V}$ 是 $\mathcal{O}_X$-模。则可定义 $\mathcal{V}[Y,p]$ 为所有被 $\mathcal{J}^p$ 的所有元素零化的截面构成的子层。于是 $\mathcal{V}[Y,p]$ 是 $\mathcal{V}$ 的 $\mathcal{O}_X$-子模。显然 $\mathcal{V}[Y,p]$ 的截面支集在 $Y$ 中，因此它们是 $\mathcal{V}$ 中所有支集在 $Y$ 中的局部截面构成的层 $\Gamma_Y(\mathcal{V})$ 的子模。显然对所有 $p\in\mathbb{Z}_+$ 有 $\mathcal{V}[Y,p]\subset\mathcal{V}[Y,p+1]$。令 $\Gamma_{[Y]}(\mathcal{V})$ 为所有子层 $\mathcal{V}[Y,p]$（$p\in\mathbb{Z}_+$）的并。则 $\Gamma_{[Y]}(\mathcal{V})$ 是 $\Gamma_Y(\mathcal{V})$ 的 $\mathcal{O}_X$-子模。

\begin{lemma}\label{lem:4.8.8}
设 $\mathcal{W}$ 是右 $\mathcal{D}_X$-模。则
$$\Gamma_{[Y]}(\mathcal{W})=i_+(i^!(\mathcal{W})).$$
\end{lemma}

\begin{proof}
设 $y\in Y$、$v\in i_+(i^!(\mathcal{W}))_y$。断言对任意 $p\in\mathbb{N}$ 有
$$\{v\in i_+(i^!(\mathcal{W}))_y\mid\operatorname{ord}(v)<p\}=\mathcal{W}[Y,p]_y.$$
显然 $\operatorname{ord}(v)=0$ 当且仅当 $v\in i^!(\mathcal{W})_y=\mathcal{W}[Y,1]_y$。这证明了 $p=1$ 的情形。设对某个 $p\ge1$ 成立。

先证明
$$\{v\in i_+(i^!(\mathcal{W}))_y\mid\operatorname{ord}(v)<p+1\}\subset\mathcal{W}[Y,p+1]_y.$$
若 $\operatorname{ord}(v)=p+1$，由引理 \ref{lem:4.8.5}(i)，对任意 $m+1\le j\le n$ 有 $\operatorname{ord}(vf_j)\le p$。故由归纳假设，对任意 $m+1\le j\le n$ 有 $vf_j\in\mathcal{W}[Y,p]_y$。由此得出 $v\in\mathcal{W}[Y,p+1]_y$。

反之，若 $v\in\mathcal{W}[Y,p+1]_y$，则对每个 $m+1\le j\le n$ 有 $vf_j\in\mathcal{W}[Y,p]_y$。由归纳假设，$vf_j\in i_+(i^!(\mathcal{W}))_y$ 且 $\operatorname{ord}(vf_j)<p$。因此
$vE=\sum_{j=m+1}^nvf_jD^j\in i_+(i^!(\mathcal{W}))_y$ 且 $\operatorname{ord}(vE)\le p$。此外，由引理 \ref{lem:4.8.5}(ii)，
$$v(E-p)f_j=vEf_j-pvf_j=v[E,f_j]+vf_jE-pvf_j=vf_j(E-(p-1))$$
是阶 $<p-1$ 的元素。故对任意 $m+1\le j\le n$ 有 $v(E-p)f_j\in\mathcal{W}[Y,p-1]_y$。因此 $v(E-p)\in\mathcal{W}[Y,p]_y$，由归纳假设 $v(E-p)\in i_+(i^!(\mathcal{W}))_y$ 且 $\operatorname{ord}(v(E-p))<p$。故 $pv=vE-v(E-p)\in i_+(i^!(\mathcal{W}))_y$ 且阶 $\le p$。由于 $p>0$，得 $v\in i_+(i^!(\mathcal{W}))_y$ 且 $\operatorname{ord}(v)<p+1$。
\end{proof}

因此 $\Gamma_{[Y]}(\mathcal{W})$ 是 $\mathcal{W}$ 的 $\mathcal{D}_X$-子模。此外，由于 $i_+$ 是正合函子而 $i^!$ 是左正合的，立即得到以下推论。

\begin{proposition}\label{prop:4.8.9}
函子 $\Gamma_{[Y]}:\mathcal{M}(\mathcal{D}_X)\to\mathcal{M}(\mathcal{D}_X)$ 是左正合的。
\end{proposition}

最后有如下结果，它给出了两个伴随态射的描述。

\begin{theorem}\label{thm:4.8.10}
设 $Y$ 是光滑簇 $X$ 的闭光滑子簇，$i:Y\to X$ 是自然包含。则
$i^!\circ i_+\cong\operatorname{id}$ 且 $i_+\circ i^!\cong\Gamma_{[Y]}$。
\end{theorem}

另一方面，有如下简单事实。

\begin{lemma}\label{lem:4.8.11}
设 $\mathcal{V}$ 是拟凝聚 $\mathcal{O}_X$-模。则 $\Gamma_{[Y]}(\mathcal{V})=\Gamma_Y(\mathcal{V})$。
\end{lemma}

\begin{proof}
可设 $X$ 是仿射簇且 $Y\ne X$。设 $v\in\Gamma(X,\mathcal{V})$ 是支集在 $Y$ 中的截面。设 $g\in R(X)$ 是在 $Y$ 上消没的非零函数，$U$ 是 $g$ 对应的主开集。则 $v|_U=0$。又 $\Gamma(U,\mathcal{V})=\Gamma(X,\mathcal{V})_g$，故对充分大的 $p\in\mathbb{Z}_+$ 有 $g^pv=0$。
\end{proof}

因此，若把最后结果应用于拟凝聚 $\mathcal{D}_X$-模，得到以下结果。

\begin{proposition}\label{prop:4.8.12}
若 $\mathcal{V}$ 是拟凝聚右 $\mathcal{D}_X$-模，则 $\Gamma_Y(\mathcal{V})=i_+(i^!(\mathcal{V}))$。
\end{proposition}

特别地，若 $\mathcal{V}$ 是支集包含在 $Y$ 中的拟凝聚右 $\mathcal{D}_X$-模，则 $i_+(i^!(\mathcal{V}))=\mathcal{V}$。这证明了以下属于 Kashiwara 的结果。

以 $\mathcal{M}^R_{qc,Y}(\mathcal{D}_X)$（相应地 $\mathcal{M}^R_{coh,Y}(\mathcal{D}_X)$，$\operatorname{Hol}_Y(\mathcal{D}_X)$）记
$\mathcal{M}^R_{qc}(\mathcal{D}_X)$（相应地 $\mathcal{M}^R_{coh}(\mathcal{D}_X)$，$\operatorname{Hol}(\mathcal{D}_X)$）中由支集在 $Y$ 中的模构成的全子范畴。

\begin{theorem}[Kashiwara 定理]\label{thm:4.8.13}
\begin{enumerate}[label=(\roman*)]
  \item 函子 $i_+:\mathcal{M}^R_{qc}(\mathcal{D}_Y)\to\mathcal{M}^R_{qc}(\mathcal{D}_X)$ 是
    $\mathcal{M}^R_{qc}(\mathcal{D}_Y)$ 与 $\mathcal{M}^R_{qc,Y}(\mathcal{D}_X)$ 之间范畴的等价。函子 $i^!$ 是 $i_+$ 的拟逆。
  \item 这些等价诱导出 $\mathcal{M}^R_{coh}(\mathcal{D}_Y)$ 与 $\mathcal{M}^R_{coh,Y}(\mathcal{D}_X)$ 之间保持完整亏格的等价。
\end{enumerate}
\end{theorem}

\begin{proof}
(i) 立即由命题 \ref{prop:4.8.3}、定理 \ref{thm:4.8.10} 与命题 \ref{prop:4.8.12} 得出。
(ii) 由前述结果得出。
\end{proof}

\section{D-模的局部上同调}

设 $Y$ 是光滑簇 $X$ 的闭光滑子簇。以 $i:Y\to X$ 记自然浸入。令 $U=X-Y$，$j:U\to X$ 是相应的开浸入。则对任意下方有界的右 $\mathcal{D}_X$-模复形 $\mathcal{V}^\cdot$，在
$D^+(\mathcal{M}_R(\mathcal{D}_X))$ 中有一个可区别三角
$$\begin{tikzcd}
& Rj_\bullet(\mathcal{V}^\cdot|_U) \arrow[dl,"+1"'] & \\
R\Gamma_Y(\mathcal{V}^\cdot) \arrow[rr] && \mathcal{V}^\cdot \arrow[lu]
\end{tikzcd}$$
且这个三角关于 $\mathcal{V}^\cdot$ 是函子性的（例如见 [3]）。

\begin{lemma}\label{lem:4.9.1}
函子 $\Gamma_Y$ 的右同调维数 $\le\dim X+1$。
\end{lemma}

\begin{proof}
由 $\mathcal{D}(\mathcal{V})$ 的上述可区别三角附带的长正合序列可知，存在从
$R^pj_\bullet(\mathcal{V})$ 到 $H^{p+1}_Y(\mathcal{V})$ 的自然满射。由于 $j_\bullet$ 的右同调维数 $\le\dim X$，结论成立。
\end{proof}

因此存在右导出函子
$R\Gamma_Y:D(\mathcal{M}_R(\mathcal{D}_X))\to D(\mathcal{M}_R(\mathcal{D}_X))$。由于
$D^+(\mathcal{M}^R_{qc}(\mathcal{D}_X))$ 等价于 $D^+(\mathcal{M}_R(\mathcal{D}_X))$ 中由具有拟凝聚上同调的复形构成的全子范畴，且内射拟凝聚 $\mathcal{D}_X$-模是松的、松层是 $\Gamma_Y$-无循环的，可见 $R\Gamma_Y$ 诱导 $D^+(\mathcal{M}^R_{qc}(\mathcal{D}_X))$ 到自身的一个正合函子，同构于
$\Gamma_Y:\mathcal{M}^R_{qc}(\mathcal{D}_X)\to\mathcal{M}^R_{qc}(\mathcal{D}_X)$ 的右导出函子。

如前，考虑伴随函子对
$i_+:\mathcal{M}^R_{qc}(\mathcal{D}_Y)\to\mathcal{M}^R_{qc}(\mathcal{D}_X)$ 与
$i^!:\mathcal{M}^R_{qc}(\mathcal{D}_X)\to\mathcal{M}^R_{qc}(\mathcal{D}_Y)$。如前所述，$i_+$ 正合而 $i^!$ 左正合。因此可考虑其导出函子
$Ri^!:D^+(\mathcal{M}^R_{qc}(\mathcal{D}_X))\to D^+(\mathcal{M}^R_{qc}(\mathcal{D}_Y))$。由前述立即得出
$$R\Gamma_Y\cong i_+\circ Ri^!$$
作为从 $D^+(\mathcal{M}^R_{qc}(\mathcal{D}_X))$ 到自身的函子。特别地，若 $\mathcal{V}$ 是拟凝聚右 $\mathcal{D}_X$-模，则对任意 $p\in\mathbb{Z}_+$ 有
$$H_Y^p(\mathcal{V})=i_+(R^pi^!(\mathcal{V})).$$
这意味着对任意 $p\in\mathbb{Z}_+$ 有
$$i^!(H^p_Y(\mathcal{V}))=i^!(i_+(R^pi^!(\mathcal{V})))=R^pi^!(\mathcal{V}),$$
这结合前述证明了以下结果。

\begin{lemma}\label{lem:4.9.2}
函子 $i^!:\mathcal{M}^R_{qc}(\mathcal{D}_X)\to\mathcal{M}^R_{qc}(\mathcal{D}_Y)$ 的右同调维数 $\le\dim X+1$。
\end{lemma}

因此存在右导出函子
$Ri^!:D(\mathcal{M}^R_{qc}(\mathcal{D}_X))\to D(\mathcal{M}^R_{qc}(\mathcal{D}_Y))$。由前述讨论得到以下结果。

\begin{theorem}\label{thm:4.9.3}
函子 $R\Gamma_Y$ 与 $i_+\circ Ri^!$ 作为从三角范畴 $D(\mathcal{M}^R_{qc}(\mathcal{D}_X))$ 到自身的正合函子是同构的。
\end{theorem}

\begin{lemma}\label{lem:4.9.4}
设 $\mathcal{V}$ 是支集在 $Y$ 中的右 $\mathcal{D}_X$-模。则 $\mathcal{V}$ 是 $\Gamma_Y$-无循环的。
\end{lemma}

\begin{proof}
如前述，由可区别三角
$$\begin{tikzcd}
& Rj_\bullet(\mathcal{D}(\mathcal{V}|_U)) \arrow[dl,"+1"'] & \\
R\Gamma_Y(\mathcal{D}(\mathcal{V})) \arrow[rr] && \mathcal{D}(\mathcal{V}) \arrow[lu]
\end{tikzcd}$$
附带的长正合序列可得，对 $p\in\mathbb{Z}_+$，$H^{p+1}_Y(\mathcal{V})$ 是
$R^pj_\bullet(\mathcal{V}|_U)$ 的商。由于 $\mathcal{V}|_U=0$，这意味着对 $p\ge1$ 有 $H^p_Y(\mathcal{V})=0$。
\end{proof}

由于对拟凝聚右 $\mathcal{D}_X$-模 $\mathcal{V}$ 有 $R^pi^!(\mathcal{V})=i^!(H_Y^p(\mathcal{V}))$（$p\in\mathbb{Z}_+$），这立即蕴含以下结果。

\begin{corollary}\label{cor:4.9.5}
设 $\mathcal{V}$ 是支集在 $Y$ 中的拟凝聚右 $\mathcal{D}_X$-模。则 $\mathcal{V}$ 是 $i^!$-无循环的。
\end{corollary}

如前，以 $\mathcal{M}^R_{qc,Y}(\mathcal{D}_X)$ 记 $\mathcal{M}^R_{qc}(\mathcal{D}_X)$ 中由支集包含在 $Y$ 中的模构成的全子范畴。这个范畴是 $\mathcal{M}^R_{qc}(\mathcal{D}_X)$ 的厚子范畴。以
$D_Y^*(\mathcal{M}^R_{qc}(\mathcal{D}_X))$ 记 $D^*(\mathcal{M}^R_{qc}(\mathcal{D}_X))$ 中由满足所有 $p\in\mathbb{Z}$ 时 $H^p(\mathcal{V}^\cdot)$ 支集在 $Y$ 中的复形 $\mathcal{V}^\cdot$ 构成的全子范畴。由前述，这是 $D^*(\mathcal{M}^R_{qc}(\mathcal{D}_X))$ 的三角子范畴。显然对任意拟凝聚右 $\mathcal{D}_X$-模复形 $\mathcal{V}^\cdot$，$R\Gamma_Y(\mathcal{V}^\cdot)$ 在
$D_Y(\mathcal{M}^R_{qc}(\mathcal{D}_X))$ 中。

现考虑 $R\Gamma_Y$ 到 $D(\mathcal{M}^R_{qc}(\mathcal{D}_X))$ 上恒等函子的自然变换。它在 $D_Y(\mathcal{M}^R_{qc}(\mathcal{D}_X))$ 上的限制诱导
$R\Gamma_Y:D_Y(\mathcal{M}^R_{qc}(\mathcal{D}_X))\to D_Y(\mathcal{M}^R_{qc}(\mathcal{D}_X))$ 到恒等函子的自然变换 $\alpha$。

\begin{proposition}\label{prop:4.9.6}
自然变换 $\alpha:R\Gamma_Y:D_Y(\mathcal{M}^R_{qc}(\mathcal{D}_X))\to D_Y(\mathcal{M}^R_{qc}(\mathcal{D}_X))$ 到恒等函子是函子的同构。
\end{proposition}

\begin{proof}
先设 $\mathcal{V}^\cdot\in D_Y(\mathcal{M}^R_{qc}(\mathcal{D}_X))$ 是下方有界复形。则由可区别三角
$$\begin{tikzcd}
& Rj_\bullet(\mathcal{V}^\cdot|_U) \arrow[dl,"+1"'] & \\
R\Gamma_Y(\mathcal{V}^\cdot) \arrow[rr] && \mathcal{V}^\cdot \arrow[lu]
\end{tikzcd}$$
可见该陈述等价于 $Rj_\bullet(\mathcal{V}^\cdot|_U)=0$。但这是显然的，因为对所有 $p\in\mathbb{Z}$ 有
$H^p(\mathcal{V}^\cdot|_U)=H^p(\mathcal{V}^\cdot)|_U=0$。一般情形由前述得出。
\end{proof}

这个结果结合 Kashiwara 定理，有以下直接推论。

\begin{theorem}\label{thm:4.9.7}
函子 $i_+:D(\mathcal{M}^R_{qc}(\mathcal{D}_Y))\to D_Y(\mathcal{M}^R_{qc}(\mathcal{D}_X))$ 是三角范畴的等价。拟逆是函子
$Ri^!:D_Y(\mathcal{M}^R_{qc}(\mathcal{D}_X))\to D(\mathcal{M}^R_{qc}(\mathcal{D}_Y))$。
\end{theorem}

\begin{proof}
由定理 \ref{thm:4.9.3} 与命题 \ref{prop:4.9.6} 可知 $i_+\circ Ri^!$ 同构于 $D_Y(\mathcal{M}^R_{qc}(\mathcal{D}_X))$ 上的恒等函子。另一方面，由于
$i_+(\mathcal{W})$（$\mathcal{W}\in\mathcal{M}^R_{qc}(\mathcal{D}_Y)$）形式的模由推论 \ref{cor:4.9.5} 是 $i^!$-无循环的，也得出 $Ri^!\circ i_+$ 同构于恒等函子。
\end{proof}

特别地，$D_Y(\mathcal{M}^R_{qc}(\mathcal{D}_X))$ 中每个复形 $\mathcal{V}^\cdot$ 都同构于由支集在 $Y$ 中的 $\mathcal{D}_X$-模构成的 $i_+(Ri^!(\mathcal{V}^\cdot))$。

现考虑 $Y$ 的闭光滑子簇 $Z$ 与自然浸入 $i_1:Z\to Y$。则有以下结果。

\begin{theorem}\label{thm:4.9.8}
$R(i\circ i_1)^!\cong Ri_1^!\circ Ri^!$ 作为从 $D(\mathcal{M}^R_{qc}(\mathcal{D}_X))$ 到
$D(\mathcal{M}^R_{qc}(\mathcal{D}_Z))$ 的函子。
\end{theorem}

\begin{proof}
由前述知 $(i\circ i_1)^!=i_1^!\circ i^!$。由于 $i^!$ 是正合函子 $i_+$ 的右伴随，它把内射拟凝聚右
$\mathcal{D}_X$-模映为内射拟凝聚右 $\mathcal{D}_Y$-模。
\end{proof}

现在想证明对 $\Gamma_Y:\mathcal{M}^R_{qc}(\mathcal{D}_X)\to\mathcal{M}^R_{qc}(\mathcal{D}_X)$ 与
$i^!:\mathcal{M}^R_{qc}(\mathcal{D}_X)\to\mathcal{M}^R_{qc}(\mathcal{D}_Y)$ 的右同调维数的更精确估计。

\begin{theorem}\label{thm:4.9.9}
\begin{enumerate}[label=(\roman*)]
  \item 函子 $\Gamma_Y:\mathcal{M}^R_{qc}(\mathcal{D}_X)\to\mathcal{M}^R_{qc}(\mathcal{D}_X)$ 的右同调维数 $\le\dim X-\dim Y$；
  \item 函子 $i^!:\mathcal{M}^R_{qc}(\mathcal{D}_X)\to\mathcal{M}^R_{qc}(\mathcal{D}_Y)$ 的右同调维数 $\le\dim X-\dim Y$。
\end{enumerate}
\end{theorem}

\begin{proof}
由于对任意拟凝聚右 $\mathcal{D}_X$-模 $\mathcal{V}$ 与 $p\in\mathbb{Z}_+$ 有
$R^pi^!(\mathcal{V})=i^!(H_Y^p(\mathcal{V}))$，这些结果等价。

此外，由于层的内射性是局部性质，第一个断言显然是局部的。因此为证明它，可设 $X$ 是仿射的，带坐标系 $(f_1,\dots,f_n;D_1,\dots,D_n)$，且 $Y$ 是 $f_{m+1},\dots,f_n$ 的公共零点集。

对 $n-m$ 作归纳。先考虑 $m=n-1$ 的情形。此时
$U=\{x\in X\mid f_n(x)\ne0\}$，即 $X$ 中主开集。故 $U$ 仿射且 $j:U\to X$ 是仿射态射。这意味着对任意拟凝聚右 $\mathcal{D}_X$-模 $\mathcal{V}$ 与 $p\ge1$ 有
$R^pj_\bullet(\mathcal{V}|_U)=0$。由于对 $p\in\mathbb{Z}_+$，$H^{p+1}_Y(\mathcal{V})$ 是
$R^pj_\bullet(\mathcal{V}|_U)$ 的商，可见对 $p\ge2$ 有 $H_Y^p(\mathcal{V})=0$。因此 (i)、(ii) 在此情形成立。

现设 $m<n-1$。则可考虑光滑子簇 $Z=\{x\in X\mid f_n(x)=0\}$。设 $i_1:Y\to Z$、$i_2:Z\to X$ 是典范浸入。由归纳假设，$i_1^!$、$i_2^!$ 具有性质 (ii)，故由定理 \ref{thm:4.9.8}，它们的复合 $i$ 也具有性质 (ii)。
\end{proof}

最后指出基变换的一个特殊情形。

\begin{lemma}\label{lem:4.9.10}
设 $Y$ 是 $X$ 的闭光滑子簇，$U=X-Y$，$i:Y\to X$、$j:U\to X$ 是自然浸入。则在
$D(\mathcal{M}_{qc}(\mathcal{D}_U))$ 上 $Ri^!\circ Rj_\bullet=0$。
\end{lemma}

\begin{proof}
设 $\mathcal{W}^\cdot$ 是下方有界的拟凝聚 $\mathcal{D}_U$-模复形。则存在 $j_\bullet$-无循环复形 $\mathcal{J}^\cdot$ 与拟同构
$\mathcal{W}^\cdot\to\mathcal{J}^\cdot$。故有交换图
$$\begin{tikzcd}
Rj_\bullet(\mathcal{W}^\cdot) \arrow[r] \arrow[d,"\sim"'] &
Rj_\bullet(Rj_\bullet(\mathcal{W}^\cdot)|_U) \arrow[d,"\sim"]\\
j_\bullet(\mathcal{J}^\cdot) \arrow[r,"="'] & j_\bullet(j_\bullet(\mathcal{J}^\cdot)|_U)
\end{tikzcd}$$
其中竖线是同构，下方水平箭头是恒等。这说明上方水平箭头也是同构。由可区别三角
$$\begin{tikzcd}
& Rj_\bullet(Rj_\bullet(\mathcal{W}^\cdot)|_U) \arrow[dl,"+1"'] & \\
R\Gamma_Y(Rj_\bullet(\mathcal{W}^\cdot)) \arrow[rr] && Rj_\bullet(\mathcal{W}^\cdot) \arrow[lu]
\end{tikzcd}$$
得出 $R\Gamma_Y(Rj_\bullet(\mathcal{W}^\cdot))=0$。于是由前述，在
$D^+(\mathcal{M}^R_{qc}(\mathcal{D}_U))$ 上 $i_+\circ Ri^!\circ Rj_\bullet=0$。再由定理 \ref{thm:4.9.7}，
$$0=Ri^!\circ i_+\circ Ri^!\circ Rj_\bullet\cong Ri^!\circ Rj_\bullet$$
在 $D^+(\mathcal{M}^R_{qc}(\mathcal{D}_U))$ 上。由标准截断点论证得出在 $D(\mathcal{M}^R_{qc}(\mathcal{D}_U))$ 上同样成立。
\end{proof}
\begin{annote}
$R\Gamma_Z$ 是“支集落在 $Z$”的截面的导出函子；在 D-模语境里它是 $R\mathcal H om_{\mathcal D_X}(\mathcal D_{X,\infty Z},-)$ 的另一种写法，$Z$ 闭、$\mathcal D_{X,\infty Z}$ 为 $Z$ 上无穷阶极点的算子层。最容易混淆的是它与 $\mathcal D_X$-模的 $\mathcal H^0_Z$（代数局部上同调）之间的换位：习题式的等式 $R\Gamma_Z(\mathcal M)\cong R\mathcal H om(\mathcal O_{X,\infty Z},\mathcal M)$ 只在 $\mathcal M$ 拟凝聚时成立。Stacks 的对应条目是 dualizing complex 一章里 torsion 与 complete 的对偶关系，可作形式对照。[来源：Stacks Project Tag 0A6Y（Lemma 47.12.3, torsion vs. complete）；原书 §9, p127–130]
\end{annote}



\section{基变换}

设 $X$、$Y$、$S$ 是三个代数簇，$\varphi:X\to S$、$\psi:Y\to S$ 是簇的两个态射。则可考虑如下代数簇的交换图
$$\begin{tikzcd}
X\times Y \arrow[r,"p_1"] \arrow[d,"p_2"' ] & X \arrow[d,"\varphi"]\\
Y \arrow[r,"\psi"'] & S
\end{tikzcd}$$
其中 $p_1$、$p_2$ 分别是到第一、第二因子的投影。设 $\Delta_S$ 是 $S\times S$ 中的对角线。由于 $S$ 是簇，$\Delta_S$ 是 $S\times S$ 的闭子簇。令
$$X\times_SY=\{(x,y)\in X\times Y\mid\varphi(x)=\psi(y)\}$$
$$=\{(x,y)\in X\times Y\mid(\varphi\circ p_1)(x,y)=(\psi\circ p_2)(x,y)\}
=((\varphi\circ p_1)\times(\psi\circ p_2))^{-1}(\Delta_S),$$
则这是 $X\times Y$ 的闭子簇。称 $X\times_SY$ 为 $X$ 与 $Y$ 在 $S$ 上的\emph{纤维积}。投影 $p_1$、$p_2$ 诱导纤维积 $X\times_SY$ 到 $X$、$Y$ 的态射，使得交换图
$$\begin{tikzcd}
X\times_SY \arrow[r,"\psi'"] \arrow[d,"\varphi'"' ] & X \arrow[d,"\varphi"]\\
Y \arrow[r,"\psi"'] & S
\end{tikzcd}$$
交换。纤维积具有以下泛性质。设 $T$ 是代数簇，$\alpha:T\to X$、$\beta:T\to Y$ 是使得交换图
$$\begin{tikzcd}
T \arrow[r,"\alpha"] \arrow[d,"\beta"' ] & X \arrow[d,"\varphi"]\\
Y \arrow[r,"\psi"'] & S
\end{tikzcd}$$
交换的簇的态射。则由积的泛性质得到一个态射 $\gamma:T\to X\times Y$ 使得 $\alpha=p_1\circ\gamma$、$\beta=p_2\circ\gamma$。因此 $\gamma(T)\subset X\times_SY$，$\gamma$ 诱导态射 $\delta:T\to X\times_SY$ 使得 $\alpha=\psi'\circ\delta$、$\beta=\varphi'\circ\delta$。

设 $Z$ 是另一个代数簇，$\chi:Z\to Y$。则有如下交换图
$$\begin{tikzcd}
(X\times_SY)\times_YZ \arrow[r,"\chi'"] \arrow[d,"\varphi''"'] &
X\times_SY \arrow[r,"\psi'"] \arrow[d,"\varphi'"'] & X \arrow[d,"\varphi"]\\
Z \arrow[r,"\chi"'] & Y \arrow[r,"\psi"'] & S
\end{tikzcd}$$
由纤维积的泛性质，得到态射 $\omega:(X\times_SY)\times_YZ\to X\times_SZ$。

\begin{lemma}\label{lem:4.10.1}
态射 $\omega:(X\times_SY)\times_YZ\to X\times_SZ$ 是同构。
\end{lemma}

\begin{proof}
由上述讨论，
$$(X\times_SY)\times_YZ=\{(x,y,z)\in X\times Y\times Z\mid\varphi'(x,y)=\chi(z),\ (x,y)\in X\times_SY\}$$
$$=\{(x,y,z)\in X\times Y\times Z\mid y=\chi(z),\ \varphi(x)=\psi(y)\}$$
且 $\omega(x,y,z)=(x,z)$。因此 $\omega$ 的像等于
$$X\times_SZ=\{(x,z)\in X\times Z\mid\varphi(x)=\psi(\chi(z))\},$$
其逆映射由 $(p_1,\chi\circ p_2,p_2)$ 在 $X\times_SY$ 上的限制给出。
\end{proof}

现考虑两种特殊情形。设 $Y$ 是 $S$ 的闭子簇，$i:Y\to S$ 是自然浸入。则
$$X\times_SY=\{(x,y)\in X\times Y\mid\varphi(x)=y\}\subset X\times S$$
等于 $\varphi$ 的图 $\Gamma_\varphi$ 与 $X\times Y$ 的交。因此 $\psi':X\times_SY\to X$ 是到 $X$ 的闭子簇 $\varphi^{-1}(Y)$ 上的同构。故纤维积图如下
$$\begin{tikzcd}
\varphi^{-1}(Y) \arrow[r,"i'"] \arrow[d,"\varphi'"' ] & X \arrow[d,"\varphi"]\\
Y \arrow[r,"i"'] & S
\end{tikzcd}$$
其中 $\varphi'$ 是 $\varphi$ 的限制，$i'$ 是自然包含。

设 $Y=Z\times S$，$\psi:Y\to S$ 是到第二个变量的投影。则可把 $Y$ 看作 $\alpha:Z\to\{pt\}$ 与 $\beta:S\to\{pt\}$ 的纤维积。故由引理 \ref{lem:4.10.1}，
$$(Z\times S)\times_SX=Z\times X,$$
且纤维积图如下
$$\begin{tikzcd}
Z\times X \arrow[r,"p_2"] \arrow[d,"\operatorname{id}\times\varphi"'] & X \arrow[d,"\varphi"]\\
Z\times S \arrow[r,"\psi"'] & S
\end{tikzcd}$$

若 $\psi:Y\to S$ 是任意态射，可用图分解 $\psi=\alpha\circ\beta$，其中 $\beta$ 是 $Y$ 到图 $\Gamma_\varphi\subset Y\times S$ 的同构，$\alpha:Y\times S\to S$ 是到第二因子的投影。故由引理 \ref{lem:4.10.1}，每个纤维积都可看作分两步得到，每步都是上述特殊类型之一。

本节的主要结果如下定理。

\begin{theorem}\label{thm:4.10.2}
设 $X$、$Y$、$S$ 是光滑代数簇，$\varphi:X\to S$、$\psi:Y\to S$ 是代数簇的态射，且纤维积 $X\times_SY$ 是光滑代数簇。则交换图
$$\begin{tikzcd}
X\times_SY \arrow[r,"\psi'"] \arrow[d,"\varphi'"' ] & X \arrow[d,"\varphi"]\\
Y \arrow[r,"\psi"'] & S
\end{tikzcd}$$
决定了从 $D(\mathcal{M}_{qc}(\mathcal{D}_X))$ 到 $D(\mathcal{M}_{qc}(\mathcal{D}_Y))$ 的函子的同构
$$\psi^!\circ\varphi_+\cong\varphi'_+\circ\psi'^!.$$
\end{theorem}

\begin{proof}
如前述，纤维积的构造总可分两步：第一步一个态射是投影，第二步是闭浸入。在积的情形下，纤维积的光滑性是自动的；在闭浸入的情形下，这是一个附加条件。

先考虑后一种情形。设 $i:Y\to S$ 是闭浸入，且 $\varphi^{-1}(Y)$ 是 $X$ 的光滑闭子簇。设 $i'$ 是 $\varphi^{-1}(Y)$ 到 $X$ 的自然浸入。令 $U=S-Y$，$j:U\to S$ 是自然浸入。又有 $X-\varphi^{-1}(Y)=\varphi^{-1}(U)$。则有如下交换图
$$\begin{tikzcd}
\varphi^{-1}(Y) \arrow[r,"i'"] & X & \arrow[l,"j'"'] \varphi^{-1}(U) \arrow[d,"\varphi''"]\\
Y \arrow[r,"i"'] \arrow[u,"\varphi'"'] & S & \arrow[l,"j"'] U
\end{tikzcd}$$
其中 $\varphi'$、$\varphi''$ 是 $\varphi$ 的相应限制。

先设 $\mathcal{V}^\cdot$ 是下方有界的拟凝聚 $\mathcal{D}_X$-模复形。考虑 $D(\mathcal{M}_{qc}(\mathcal{D}_X))$ 中的可区别三角
$$\begin{tikzcd}
& Rj'_\bullet(\mathcal{V}^\cdot|_{\varphi^{-1}(U)}) \arrow[dl,"+1"'] & \\
R\Gamma_{\varphi^{-1}(Y)}(\mathcal{V}^\cdot) \arrow[rr] && \mathcal{V}^\cdot \arrow[lu]
\end{tikzcd}$$
应用导出函子 $\varphi_+$，得到 $D(\mathcal{M}_{qc}(\mathcal{D}_S))$ 中的可区别三角
$$\begin{tikzcd}
& \varphi_+(Rj'_\bullet(\mathcal{V}^\cdot|_{\varphi^{-1}(U)})) \arrow[dl,"+1"'] & \\
\varphi_+(R\Gamma_{\varphi^{-1}(Y)}(\mathcal{V}^\cdot)) \arrow[rr] && \varphi_+(\mathcal{V}^\cdot) \arrow[lu]
\end{tikzcd}$$
再应用导出函子 $Ri^!$，得到 $D(\mathcal{M}_{qc}(\mathcal{D}_Y))$ 中的可区别三角
$$\begin{tikzcd}
& Ri^!(\varphi_+(Rj'_\bullet(\mathcal{V}^\cdot|_{\varphi^{-1}(U)}))) \arrow[dl,"+1"'] & \\
Ri^!(\varphi_+(R\Gamma_{\varphi^{-1}(Y)}(\mathcal{V}^\cdot))) \arrow[rr] && Ri^!(\varphi_+(\mathcal{V}^\cdot)) \arrow[lu]
\end{tikzcd}$$
由于
$$\varphi_+(Rj'_\bullet(\mathcal{V}^\cdot|_{\varphi^{-1}(U)}))
=(\varphi\circ j')_+(\mathcal{V}^\cdot|_{\varphi^{-1}(U)})
=(j\circ\varphi'')_+(\mathcal{V}^\cdot|_{\varphi^{-1}(U)})$$
$$=Rj_\bullet(\varphi''_+(\mathcal{V}^\cdot|_{\varphi^{-1}(U)}))=Rj_\bullet(\varphi_+(\mathcal{V}^\cdot)|_U),$$
可见
$$Ri^!(\varphi_+(Rj'_\bullet(\mathcal{V}^\cdot|_{\varphi^{-1}(U)})))=Ri^!(Rj_\bullet(\varphi_+(\mathcal{V}^\cdot)|_U)).$$
由引理 \ref{lem:4.9.10}，这等于 0。因此自然态射
$$Ri^!(\varphi_+(R\Gamma_{\varphi^{-1}(Y)}(\mathcal{V}^\cdot)))\longrightarrow Ri^!(\varphi_+(\mathcal{V}^\cdot))$$
是关于 $\mathcal{V}^\cdot$ 函子性的同构。由于所涉函子都具有有限同调维数，由截断点论证可知这个自然态射对任意拟凝聚 $\mathcal{D}_X$-模复形 $\mathcal{V}^\cdot$ 都是同构。

另一方面，由前述有
$$Ri^!(\varphi_+(R\Gamma_{\varphi^{-1}(Y)}(\mathcal{V}^\cdot)))
=Ri^!(\varphi_+(i'_+(Ri'^!(\mathcal{V}^\cdot))))
=Ri^!((\varphi\circ i')_+(Ri'^!(\mathcal{V}^\cdot)))$$
$$=Ri^!((i\circ\varphi')_+(Ri'^!(\mathcal{V}^\cdot)))
=Ri^!(i_+(\varphi'_+(Ri'^!(\mathcal{V}^\cdot))))=\varphi'_+(Ri'^!(\mathcal{V}^\cdot)).$$
因此在这个特殊情形下结论得证。

现考虑投影的情形。若 $p:Z\times S\to S$ 是到第二因子的投影，如前得到纤维积交换图
$$\begin{tikzcd}
Z\times X \arrow[r,"p_2"] \arrow[d,"\operatorname{id}_Z\times\varphi"'] & X \arrow[d,"\varphi"]\\
Z\times S \arrow[r,"p"'] & S
\end{tikzcd}$$
于是由前述有
$$(p^!\circ\varphi_+)(\mathcal{V}^\cdot)=Lp_+(\varphi_+(\mathcal{V}^\cdot))[-\dim S]
=\mathcal{O}_Z\boxtimes\varphi_+(\mathcal{V}^\cdot)[-\dim S]$$
$$=(\operatorname{id}_Z\times\varphi)_+\bigl((\mathcal{O}_Z\boxtimes\mathcal{V}^\cdot)[-\dim S]\bigr)
=(\operatorname{id}_Z\times\varphi)_+\bigl(p_2^+(\mathcal{V}^\cdot)[-\dim S]\bigr)$$
$$=(\operatorname{id}_Z\times\varphi)_+(p_2^!(\mathcal{V}^\cdot))
=((\operatorname{id}_Z\times\varphi)_+\circ p_2^!)(\mathcal{V}^\cdot),$$
这在此情形下建立了基变换。
\end{proof}

这使我们能够推广引理 \ref{lem:4.9.10}。

\begin{corollary}\label{cor:4.10.3}
设 $X$、$Y$、$Z$ 是光滑代数簇，$\varphi:X\to Z$、$\psi:Y\to Z$ 是代数簇的态射，使得
$\varphi(X)\cap\psi(Y)=\varnothing$。则在 $D(\mathcal{M}_{qc}(\mathcal{D}_X))$ 上 $\psi^!\circ\varphi_+=0$。
\end{corollary}
